\documentclass[10pt,a4paper,twocolumn]{article}
\usepackage[utf8]{inputenc}
\usepackage{geometry}
\geometry{a4paper, margin=0.4in}
\usepackage{xcolor}
\usepackage{tcolorbox}
\tcbuselibrary{skins}
\usepackage{amsmath,amssymb}
\usepackage{mathptmx}
\usepackage{float}
\setlength{\columnsep}{0.4in}
\setlength{\parindent}{0pt}
\definecolor{UnitColorOne}{rgb}{0.2, 0.6, 0.8}
\definecolor{UnitColorTwo}{rgb}{0.1, 0.65, 0.3}
\definecolor{UnitColorThree}{rgb}{0.9, 0.45, 0.1}
\definecolor{UnitColorFour}{rgb}{0.8, 0.1, 0.2}
\definecolor{UnitColorFive}{rgb}{0.6, 0.3, 0.8}
\definecolor{UnitColorSix}{rgb}{0.2, 0.5, 0.2}
\definecolor{UnitColorSeven}{rgb}{0.1, 0.3, 0.7}
\definecolor{UnitColorEight}{rgb}{0.65, 0.45, 0.2}
\definecolor{UnitColorNine}{rgb}{0.7, 0.1, 0.5}
\definecolor{UnitColorTen}{rgb}{0.4, 0.5, 0.6}
\begin{document}

\twocolumn[
  \begin{center}
    \Huge\bfseries\sffamily OE-EC803A: Internet of Things (IoT) \\
    \vskip 0.1in
    \Large\bfseries Active Recall Flashcards (Units I--X) \\
    \vskip 0.05in
    \large Semester VIII B.Tech ECE (MAKAUT) \\
    \vskip 0.2in
    \hrule
    \vskip 0.2in
  \end{center}
]

\section*{Unit I: Chapter 01 Flashcards}

\begin{tcolorbox}[
  colback=white,
  colframe=UnitColorOne,
  colbacktitle=gray!10!white,
  coltitle=black,
  fonttitle=\bfseries\sffamily\small,
  title={Unit I $\bullet$ Card 1},
  boxrule=1.2pt,
  arc=2mm,
  outer arc=2mm,
  boxsep=1.5mm,
  top=1mm,
  bottom=1mm,
  left=1.5mm,
  right=1.5mm,
  before skip=2.5mm,
  after skip=2.5mm
]
\textbf{Q:} What is the Internet of Things (IoT)?
\tcbline
\textbf{Ans:} A dynamic global network infrastructure with self-configuring capabilities based on standard and interoperable communication protocols, connecting physical and virtual "things" embedded with sensors, actuators, and connectivity.
\end{tcolorbox}


\begin{tcolorbox}[
  colback=white,
  colframe=UnitColorOne,
  colbacktitle=gray!10!white,
  coltitle=black,
  fonttitle=\bfseries\sffamily\small,
  title={Unit I $\bullet$ Card 2},
  boxrule=1.2pt,
  arc=2mm,
  outer arc=2mm,
  boxsep=1.5mm,
  top=1mm,
  bottom=1mm,
  left=1.5mm,
  right=1.5mm,
  before skip=2.5mm,
  after skip=2.5mm
]
\textbf{Q:} What is the fundamental difference between the World Wide Web (WWW) and the Internet of Things (IoT)?
\tcbline
\textbf{Ans:} The WWW connects \textbf{humans to information} (browser-based client-server requests), while the IoT connects \textbf{devices/things to other devices/things} (machine-to-machine data exchanges and actions).
\end{tcolorbox}


\begin{tcolorbox}[
  colback=white,
  colframe=UnitColorOne,
  colbacktitle=gray!10!white,
  coltitle=black,
  fonttitle=\bfseries\sffamily\small,
  title={Unit I $\bullet$ Card 3},
  boxrule=1.2pt,
  arc=2mm,
  outer arc=2mm,
  boxsep=1.5mm,
  top=1mm,
  bottom=1mm,
  left=1.5mm,
  right=1.5mm,
  before skip=2.5mm,
  after skip=2.5mm
]
\textbf{Q:} Explain the concept of an "Enchanted Object" and give an example.
\tcbline
\textbf{Ans:} A physical, everyday object enhanced with digital intelligence, sensors, and internet connectivity, providing information via subtle, ambient cues without changing its native physical form. \textbf{Example:} An umbrella handle that glows blue when rain is forecast.
\end{tcolorbox}


\begin{tcolorbox}[
  colback=white,
  colframe=UnitColorOne,
  colbacktitle=gray!10!white,
  coltitle=black,
  fonttitle=\bfseries\sffamily\small,
  title={Unit I $\bullet$ Card 4},
  boxrule=1.2pt,
  arc=2mm,
  outer arc=2mm,
  boxsep=1.5mm,
  top=1mm,
  bottom=1mm,
  left=1.5mm,
  right=1.5mm,
  before skip=2.5mm,
  after skip=2.5mm
]
\textbf{Q:} Differentiate between a Wireless Sensor Network (WSN) and an IoT network.
\tcbline
\textbf{Ans:} A WSN is a localized network of sensor nodes focused purely on data collection (sensing). An IoT network is a global network that integrates sensors, internet gateways, cloud analytics, and \textbf{actuators} to both sense and dynamically control/affect the physical world.
\end{tcolorbox}


\begin{tcolorbox}[
  colback=white,
  colframe=UnitColorOne,
  colbacktitle=gray!10!white,
  coltitle=black,
  fonttitle=\bfseries\sffamily\small,
  title={Unit I $\bullet$ Card 5},
  boxrule=1.2pt,
  arc=2mm,
  outer arc=2mm,
  boxsep=1.5mm,
  top=1mm,
  bottom=1mm,
  left=1.5mm,
  right=1.5mm,
  before skip=2.5mm,
  after skip=2.5mm
]
\textbf{Q:} Compare Machine-to-Machine (M2M) communication and the Internet of Things (IoT) across scope and protocols.
\tcbline
\textbf{Ans:} M2M is typically point-to-point, closed, isolated, and hardware-centric using legacy/cellular links. IoT is a global, open-standard, cloud-centric network using standard internet protocols (IP, TCP/IP, MQTT, CoAP) and web APIs.
\end{tcolorbox}


\begin{tcolorbox}[
  colback=white,
  colframe=UnitColorOne,
  colbacktitle=gray!10!white,
  coltitle=black,
  fonttitle=\bfseries\sffamily\small,
  title={Unit I $\bullet$ Card 6},
  boxrule=1.2pt,
  arc=2mm,
  outer arc=2mm,
  boxsep=1.5mm,
  top=1mm,
  bottom=1mm,
  left=1.5mm,
  right=1.5mm,
  before skip=2.5mm,
  after skip=2.5mm
]
\textbf{Q:} Explain the four "flavors" (communication models) of IoT.
\tcbline
\textbf{Ans:} \begin{enumerate}
  \item \textbf{Device-to-Device (D2D):} Direct local link between devices (e.g. Bluetooth).
  \item \textbf{Device-to-Cloud (D2C):} Direct IP connection to a cloud platform.
  \item \textbf{Device-to-Gateway (D2G):} Local link to a gateway, which forwards data to the cloud.
  \item \textbf{Back-End Data Sharing (C2C):} Cloud-to-cloud API connections for data sharing.
\end{enumerate}
\end{tcolorbox}


\begin{tcolorbox}[
  colback=white,
  colframe=UnitColorOne,
  colbacktitle=gray!10!white,
  coltitle=black,
  fonttitle=\bfseries\sffamily\small,
  title={Unit I $\bullet$ Card 7},
  boxrule=1.2pt,
  arc=2mm,
  outer arc=2mm,
  boxsep=1.5mm,
  top=1mm,
  bottom=1mm,
  left=1.5mm,
  right=1.5mm,
  before skip=2.5mm,
  after skip=2.5mm
]
\textbf{Q:} What are the key differences between traditional IoT and Industrial IoT (IIoT)?
\tcbline
\textbf{Ans:} Traditional IoT is consumer-focused (convenience, moderate reliability, low safety risk). IIoT is industrial-focused (mission-critical operations, zero-downtime tolerance, extreme safety/security standards, high financial impact of failure).
\end{tcolorbox}


\begin{tcolorbox}[
  colback=white,
  colframe=UnitColorOne,
  colbacktitle=gray!10!white,
  coltitle=black,
  fonttitle=\bfseries\sffamily\small,
  title={Unit I $\bullet$ Card 8},
  boxrule=1.2pt,
  arc=2mm,
  outer arc=2mm,
  boxsep=1.5mm,
  top=1mm,
  bottom=1mm,
  left=1.5mm,
  right=1.5mm,
  before skip=2.5mm,
  after skip=2.5mm
]
\textbf{Q:} Describe the role of a "Gateway" in an IoT network.
\tcbline
\textbf{Ans:} An IoT Gateway acts as a bridge between local constrained devices and the wider internet. It performs protocol translation (e.g., Zigbee to IP), data aggregation, and local preprocessing (edge computing).
\end{tcolorbox}


\begin{tcolorbox}[
  colback=white,
  colframe=UnitColorOne,
  colbacktitle=gray!10!white,
  coltitle=black,
  fonttitle=\bfseries\sffamily\small,
  title={Unit I $\bullet$ Card 9},
  boxrule=1.2pt,
  arc=2mm,
  outer arc=2mm,
  boxsep=1.5mm,
  top=1mm,
  bottom=1mm,
  left=1.5mm,
  right=1.5mm,
  before skip=2.5mm,
  after skip=2.5mm
]
\textbf{Q:} List the three interfaces defined in the ETSI M2M Functional Architecture.
\tcbline
\textbf{Ans:} \begin{enumerate}
  \item \textbf{dIa (Device-to-Application):} Connects device apps to local service capabilities.
  \item \textbf{mIa (M2M-to-Application):} Connects network apps to network service capabilities.
  \item \textbf{mId (M2M-to-Device):} Connects the device/gateway layer to the network service capability layer.
\end{enumerate}
\end{tcolorbox}


\begin{tcolorbox}[
  colback=white,
  colframe=UnitColorOne,
  colbacktitle=gray!10!white,
  coltitle=black,
  fonttitle=\bfseries\sffamily\small,
  title={Unit I $\bullet$ Card 10},
  boxrule=1.2pt,
  arc=2mm,
  outer arc=2mm,
  boxsep=1.5mm,
  top=1mm,
  bottom=1mm,
  left=1.5mm,
  right=1.5mm,
  before skip=2.5mm,
  after skip=2.5mm
]
\textbf{Q:} What are the three core functional layers of a Generic M2M System Solution?
\tcbline
\textbf{Ans:} \begin{enumerate}
  \item \textbf{M2M Area Network (Device Domain):} Local sensors, actuators, and local links.
  \item \textbf{Communication Network (Network Domain):} Long-distance links (Cellular, WAN, satellite).
  \item \textbf{M2M Applications (Application Domain):} Cloud backend and user applications.
\end{enumerate}
\end{tcolorbox}


\begin{tcolorbox}[
  colback=white,
  colframe=UnitColorOne,
  colbacktitle=gray!10!white,
  coltitle=black,
  fonttitle=\bfseries\sffamily\small,
  title={Unit I $\bullet$ Card 11},
  boxrule=1.2pt,
  arc=2mm,
  outer arc=2mm,
  boxsep=1.5mm,
  top=1mm,
  bottom=1mm,
  left=1.5mm,
  right=1.5mm,
  before skip=2.5mm,
  after skip=2.5mm
]
\textbf{Q:} List the four main stages of the Information-Driven Value Chain for IoT.
\tcbline
\textbf{Ans:} \textbf{Sense} (sensors capture states) $\rightarrow$ \textbf{Transmit} (send data over IP) $\rightarrow$ \textbf{Store/Analyze} (databases and data processing) $\rightarrow$ \textbf{Decide/Act} (actuators perform control actions).
\end{tcolorbox}


\begin{tcolorbox}[
  colback=white,
  colframe=UnitColorOne,
  colbacktitle=gray!10!white,
  coltitle=black,
  fonttitle=\bfseries\sffamily\small,
  title={Unit I $\bullet$ Card 12},
  boxrule=1.2pt,
  arc=2mm,
  outer arc=2mm,
  boxsep=1.5mm,
  top=1mm,
  bottom=1mm,
  left=1.5mm,
  right=1.5mm,
  before skip=2.5mm,
  after skip=2.5mm
]
\textbf{Q:} What are the SOS and SensorML standards in the OGC framework?
\tcbline
\textbf{Ans:} \textbf{SOS (Sensor Observation Service)} is an API standard for querying real-time sensor observations. \textbf{SensorML (Sensor Model Language)} is an XML schema describing sensor system parameters and metadata.
\end{tcolorbox}


\begin{tcolorbox}[
  colback=white,
  colframe=UnitColorOne,
  colbacktitle=gray!10!white,
  coltitle=black,
  fonttitle=\bfseries\sffamily\small,
  title={Unit I $\bullet$ Card 13},
  boxrule=1.2pt,
  arc=2mm,
  outer arc=2mm,
  boxsep=1.5mm,
  top=1mm,
  bottom=1mm,
  left=1.5mm,
  right=1.5mm,
  before skip=2.5mm,
  after skip=2.5mm
]
\textbf{Q:} Name four major challenges faced when implementing an IoT system.
\tcbline
\textbf{Ans:} \begin{enumerate}
  \item \textbf{Security \& Privacy} (vulnerability of resource-constrained nodes).
  \item \textbf{Interoperability} (communication barriers between vendors).
  \item \textbf{Data Volume Management} (processing massive unstructured streams).
  \item \textbf{Power constraints} (achieving long battery lives for remote nodes).
\end{enumerate}
\end{tcolorbox}


\begin{tcolorbox}[
  colback=white,
  colframe=UnitColorOne,
  colbacktitle=gray!10!white,
  coltitle=black,
  fonttitle=\bfseries\sffamily\small,
  title={Unit I $\bullet$ Card 14},
  boxrule=1.2pt,
  arc=2mm,
  outer arc=2mm,
  boxsep=1.5mm,
  top=1mm,
  bottom=1mm,
  left=1.5mm,
  right=1.5mm,
  before skip=2.5mm,
  after skip=2.5mm
]
\textbf{Q:} How does an IoT-oriented approach provide more value than M2M in a Health Band system?
\tcbline
\textbf{Ans:} M2M only sends heart rate to a single monitor. IoT integrates the data with EMR databases, third-party fitness apps, local pharmacies, and runs predictive cloud algorithms to alert the user of anomalies via their phone.
\end{tcolorbox}

\section*{Unit II: Chapter 02 Flashcards}

\begin{tcolorbox}[
  colback=white,
  colframe=UnitColorTwo,
  colbacktitle=gray!10!white,
  coltitle=black,
  fonttitle=\bfseries\sffamily\small,
  title={Unit II $\bullet$ Card 1},
  boxrule=1.2pt,
  arc=2mm,
  outer arc=2mm,
  boxsep=1.5mm,
  top=1mm,
  bottom=1mm,
  left=1.5mm,
  right=1.5mm,
  before skip=2.5mm,
  after skip=2.5mm
]
\textbf{Q:} What is the primary goal of the design principles discussed in this unit?
\tcbline
\textbf{Ans:} The primary goal is to guide the creation of "calm and ambient technology" where the technology integrates seamlessly into the user's environment, providing value without being intrusive or demanding constant attention. It emphasizes human-centered design.
\end{tcolorbox}


\begin{tcolorbox}[
  colback=white,
  colframe=UnitColorTwo,
  colbacktitle=gray!10!white,
  coltitle=black,
  fonttitle=\bfseries\sffamily\small,
  title={Unit II $\bullet$ Card 2},
  boxrule=1.2pt,
  arc=2mm,
  outer arc=2mm,
  boxsep=1.5mm,
  top=1mm,
  bottom=1mm,
  left=1.5mm,
  right=1.5mm,
  before skip=2.5mm,
  after skip=2.5mm
]
\textbf{Q:} How does the concept of "Calm Technology" differ from traditional "Ugly Technology"?
\tcbline
\textbf{Ans:} Ugly/Intrusive Technology is "noisy"—it uses active screens, requires focus, and constantly demands attention (like a constantly pinging smartphone). Calm Technology is "quiet"—it uses peripheral awareness, ambient information display, and subtle cues, allowing users to stay informed without being overwhelmed.
\end{tcolorbox}


\begin{tcolorbox}[
  colback=white,
  colframe=UnitColorTwo,
  colbacktitle=gray!10!white,
  coltitle=black,
  fonttitle=\bfseries\sffamily\small,
  title={Unit II $\bullet$ Card 3},
  boxrule=1.2pt,
  arc=2mm,
  outer arc=2mm,
  boxsep=1.5mm,
  top=1mm,
  bottom=1mm,
  left=1.5mm,
  right=1.5mm,
  before skip=2.5mm,
  after skip=2.5mm
]
\textbf{Q:} What is Ambient Technology? Give three classic examples.
\tcbline
\textbf{Ans:} Physical devices that implement Calm Technology by displaying data through subtle, environmental cues (color, light, motion) rather than pixel-based screens.  
\textbf{Examples:}  
\begin{enumerate}
  \item \textit{Ambient Orb:} A frosted ball that changes color based on stock market indices.
  \item \textit{Ambient Umbrella:} An umbrella handle that pulses blue if rain is forecast.
  \item \textit{Water Lamp:} Projects light ripples on the ceiling representing office network activity.
\end{enumerate}
\end{tcolorbox}


\begin{tcolorbox}[
  colback=white,
  colframe=UnitColorTwo,
  colbacktitle=gray!10!white,
  coltitle=black,
  fonttitle=\bfseries\sffamily\small,
  title={Unit II $\bullet$ Card 4},
  boxrule=1.2pt,
  arc=2mm,
  outer arc=2mm,
  boxsep=1.5mm,
  top=1mm,
  bottom=1mm,
  left=1.5mm,
  right=1.5mm,
  before skip=2.5mm,
  after skip=2.5mm
]
\textbf{Q:} State the five core design rules of Calm Technology.
\tcbline
\textbf{Ans:} \begin{enumerate}
  \item \textbf{Peripheral Attention Priority:} Must work in the user's peripheral sensory field.
  \item \textbf{Seamless Shift Control:} Must shift easily between the periphery and the center of focus.
  \item \textbf{Enhancing Presence:} Must provide a sense of context and presence.
  \item \textbf{Peripheral Reach Extension:} Must extend the user's awareness of remote systems.
  \item \textbf{Respecting Attention Limits:} Must work in harmony with human attention bandwidth.
\end{enumerate}
\end{tcolorbox}


\begin{tcolorbox}[
  colback=white,
  colframe=UnitColorTwo,
  colbacktitle=gray!10!white,
  coltitle=black,
  fonttitle=\bfseries\sffamily\small,
  title={Unit II $\bullet$ Card 5},
  boxrule=1.2pt,
  arc=2mm,
  outer arc=2mm,
  boxsep=1.5mm,
  top=1mm,
  bottom=1mm,
  left=1.5mm,
  right=1.5mm,
  before skip=2.5mm,
  after skip=2.5mm
]
\textbf{Q:} Detail the working principle, advantages, and limitations of David Rose's Ambient Orb.
\tcbline
\textbf{Ans:} \begin{itemize}
  \item   \textbf{Working Principle:} A frosted sphere receiving stock or traffic data feeds over the internet, changing its internal LED color accordingly (e.g., green for up, red for down).
  \item   \textbf{Advantage:} Allows instant, non-intrusive, glanceable monitoring in the background of a room.
  \item   \textbf{Limitation:} Low information density; cannot show exact numeric values.
\end{itemize}
\end{tcolorbox}


\begin{tcolorbox}[
  colback=white,
  colframe=UnitColorTwo,
  colbacktitle=gray!10!white,
  coltitle=black,
  fonttitle=\bfseries\sffamily\small,
  title={Unit II $\bullet$ Card 6},
  boxrule=1.2pt,
  arc=2mm,
  outer arc=2mm,
  boxsep=1.5mm,
  top=1mm,
  bottom=1mm,
  left=1.5mm,
  right=1.5mm,
  before skip=2.5mm,
  after skip=2.5mm
]
\textbf{Q:} Explain the concept of "Magic as Metaphor" in IoT design. Give two examples.
\tcbline
\textbf{Ans:} Using magical archetypes (magic wand, crystal ball, enchanted mirror) from myths as design metaphors to explain how invisible, wireless IoT systems operate.  
\begin{itemize}
  \item   \textbf{Examples:}
\end{itemize}
\begin{enumerate}
  \item \textit{Kymera Magic Wand:} A physical wand that uses gestures to change TV channels.
  \item \textit{Smart Mirror:} A bathroom mirror that displays digital text on its glass surface, acting like an enchanted mirror.
\end{enumerate}
\end{tcolorbox}


\begin{tcolorbox}[
  colback=white,
  colframe=UnitColorTwo,
  colbacktitle=gray!10!white,
  coltitle=black,
  fonttitle=\bfseries\sffamily\small,
  title={Unit II $\bullet$ Card 7},
  boxrule=1.2pt,
  arc=2mm,
  outer arc=2mm,
  boxsep=1.5mm,
  top=1mm,
  bottom=1mm,
  left=1.5mm,
  right=1.5mm,
  before skip=2.5mm,
  after skip=2.5mm
]
\textbf{Q:} What are the three primary critiques or limitations of using "Magic as Metaphor"?
\tcbline
\textbf{Ans:} \begin{enumerate}
  \item \textbf{Expectation Mismatch:} Users expect infinite magical capabilities, leading to disappointment.
  \item \textbf{No Clear Error States (Lack of Debuggability):} Magic does not have error codes. If a gesture fails, the user cannot easily debug if the network is down or the battery is dead.
  \item \textbf{Cultural Variance:} Magical myths vary globally, making some metaphors less intuitive in certain markets.
\end{enumerate}
\end{tcolorbox}


\begin{tcolorbox}[
  colback=white,
  colframe=UnitColorTwo,
  colbacktitle=gray!10!white,
  coltitle=black,
  fonttitle=\bfseries\sffamily\small,
  title={Unit II $\bullet$ Card 8},
  boxrule=1.2pt,
  arc=2mm,
  outer arc=2mm,
  boxsep=1.5mm,
  top=1mm,
  bottom=1mm,
  left=1.5mm,
  right=1.5mm,
  before skip=2.5mm,
  after skip=2.5mm
]
\textbf{Q:} List five distinct advantages of Calm and Ambient Technology.
\tcbline
\textbf{Ans:} \begin{enumerate}
  \item Reduces cognitive fatigue and stress.
  \item Blends naturally into the decor, avoiding screen clutter.
  \item Preserves human-to-human face-to-face interaction.
  \item Fosters ambient awareness of external systems.
  \item High energy efficiency compared to active LCD screens.
\end{enumerate}
\end{tcolorbox}


\begin{tcolorbox}[
  colback=white,
  colframe=UnitColorTwo,
  colbacktitle=gray!10!white,
  coltitle=black,
  fonttitle=\bfseries\sffamily\small,
  title={Unit II $\bullet$ Card 9},
  boxrule=1.2pt,
  arc=2mm,
  outer arc=2mm,
  boxsep=1.5mm,
  top=1mm,
  bottom=1mm,
  left=1.5mm,
  right=1.5mm,
  before skip=2.5mm,
  after skip=2.5mm
]
\textbf{Q:} List five distinct disadvantages of Calm and Ambient Technology.
\tcbline
\textbf{Ans:} \begin{enumerate}
  \item Low information density (conveys status, not detailed text).
  \item High ambiguity (abstract cues can be misread).
  \item Inappropriate for high-priority critical alerts (e.g., fire alarms).
  \item Difficult to troubleshoot and debug due to lack of screens.
  \item Accessibility issues (color blindness or hearing loss can block cues).
\end{enumerate}
\end{tcolorbox}


\begin{tcolorbox}[
  colback=white,
  colframe=UnitColorTwo,
  colbacktitle=gray!10!white,
  coltitle=black,
  fonttitle=\bfseries\sffamily\small,
  title={Unit II $\bullet$ Card 10},
  boxrule=1.2pt,
  arc=2mm,
  outer arc=2mm,
  boxsep=1.5mm,
  top=1mm,
  bottom=1mm,
  left=1.5mm,
  right=1.5mm,
  before skip=2.5mm,
  after skip=2.5mm
]
\textbf{Q:} How does "Privacy" in IoT device design extend beyond standard data encryption?
\tcbline
\textbf{Ans:} It focuses on the physical space boundaries:  
\begin{enumerate}
  \item \textbf{Consent \& Control:} Offering physical switches to disable sensor recording.
  \item \textbf{Contextual Integrity:} Ensuring household data (e.g. thermostat activity) is not sold or repurposed.
  \item \textbf{Data Minimization:} Keeping raw data collection to the minimum required.
  \item \textbf{Local Processing:} Analyzing voice/images locally at the edge/gateway rather than streaming to the cloud.
  \item \textbf{Transparency:} Using hardwired LEDs to show when microphones or cameras are actively powered.
\end{enumerate}
\end{tcolorbox}


\begin{tcolorbox}[
  colback=white,
  colframe=UnitColorTwo,
  colbacktitle=gray!10!white,
  coltitle=black,
  fonttitle=\bfseries\sffamily\small,
  title={Unit II $\bullet$ Card 11},
  boxrule=1.2pt,
  arc=2mm,
  outer arc=2mm,
  boxsep=1.5mm,
  top=1mm,
  bottom=1mm,
  left=1.5mm,
  right=1.5mm,
  before skip=2.5mm,
  after skip=2.5mm
]
\textbf{Q:} Explain "Web Thinking for Connected Devices" and name three key principles.
\tcbline
\textbf{Ans:} Designing physical IoT architectures using the open, scalable, and standardized principles of the World Wide Web.  
\begin{itemize}
  \item   \textbf{Key Principles:}
\end{itemize}
\begin{enumerate}
  \item \textit{URIs/URLs for Things:} Every sensor/actuator has a unique address.
  \item \textit{RESTful APIs:} Devices expose standard HTTP commands (GET/POST/PUT/DELETE) using JSON.
  \item \textit{Loose Coupling:} Hardware is independent of the applications controlling it.
\end{enumerate}
\end{tcolorbox}


\begin{tcolorbox}[
  colback=white,
  colframe=UnitColorTwo,
  colbacktitle=gray!10!white,
  coltitle=black,
  fonttitle=\bfseries\sffamily\small,
  title={Unit II $\bullet$ Card 12},
  boxrule=1.2pt,
  arc=2mm,
  outer arc=2mm,
  boxsep=1.5mm,
  top=1mm,
  bottom=1mm,
  left=1.5mm,
  right=1.5mm,
  before skip=2.5mm,
  after skip=2.5mm
]
\textbf{Q:} What is a "Physical Mashup" in the context of Web Thinking? Give an example.
\tcbline
\textbf{Ans:} Linking different physical devices together using standard web APIs and middleware (like IFTTT or Node-RED).  
\begin{itemize}
  \item   \textbf{Example:} "If a smart security camera detects motion, Then send an HTTP command to blink a smart light bulb."
\end{itemize}
\end{tcolorbox}


\begin{tcolorbox}[
  colback=white,
  colframe=UnitColorTwo,
  colbacktitle=gray!10!white,
  coltitle=black,
  fonttitle=\bfseries\sffamily\small,
  title={Unit II $\bullet$ Card 13},
  boxrule=1.2pt,
  arc=2mm,
  outer arc=2mm,
  boxsep=1.5mm,
  top=1mm,
  bottom=1mm,
  left=1.5mm,
  right=1.5mm,
  before skip=2.5mm,
  after skip=2.5mm
]
\textbf{Q:} Define "Affordance" and explain the "Clash of Affordances" in IoT.
\tcbline
\textbf{Ans:} \begin{itemize}
  \item   \textbf{Affordance:} The physical properties of an object that suggest how to use it (e.g., button invites pushing).
  \item   \textbf{Clash of Affordances:} A smart object looks like its traditional physical counterpart (e.g., a smart cup), but has hidden digital affordances (Bluetooth pairing, sensors) that are completely invisible.
\end{itemize}
\end{tcolorbox}


\begin{tcolorbox}[
  colback=white,
  colframe=UnitColorTwo,
  colbacktitle=gray!10!white,
  coltitle=black,
  fonttitle=\bfseries\sffamily\small,
  title={Unit II $\bullet$ Card 14},
  boxrule=1.2pt,
  arc=2mm,
  outer arc=2mm,
  boxsep=1.5mm,
  top=1mm,
  bottom=1mm,
  left=1.5mm,
  right=1.5mm,
  before skip=2.5mm,
  after skip=2.5mm
]
\textbf{Q:} What are three design solutions to resolve the Clash of Affordances?
\tcbline
\textbf{Ans:} \begin{enumerate}
  \item \textbf{Visual Signifiers:} Use subtle indicators like LED rings or textured borders to guide touch.
  \item \textbf{Sensory Feedback:} Provide immediate haptic vibes or chirps when a sensor is triggered.
  \item \textbf{Graceful Degradation:} Ensure the device still functions as its physical counterpart if the battery dies or Wi-Fi is lost (e.g. a smart lock still opens with a physical key).
\end{enumerate}
\end{tcolorbox}

\section*{Unit III: Chapter 03 Flashcards}

\begin{tcolorbox}[
  colback=white,
  colframe=UnitColorThree,
  colbacktitle=gray!10!white,
  coltitle=black,
  fonttitle=\bfseries\sffamily\small,
  title={Unit III $\bullet$ Card 1},
  boxrule=1.2pt,
  arc=2mm,
  outer arc=2mm,
  boxsep=1.5mm,
  top=1mm,
  bottom=1mm,
  left=1.5mm,
  right=1.5mm,
  before skip=2.5mm,
  after skip=2.5mm
]
\textbf{Q:} Define Bluetooth Low Energy (BLE) and name its four core features.
\tcbline
\textbf{Ans:} \textbf{BLE} is a wireless personal area network technology designed for low-power, short-range IoT communication.  
\textit{Features:} Ultra-low power consumption (runs on coin cells), fast connection setup ($<$3 ms), Advertising Mode (beacon broadcasting), and GATT data hierarchy (Profiles/Services/Characteristics).
\end{tcolorbox}


\begin{tcolorbox}[
  colback=white,
  colframe=UnitColorThree,
  colbacktitle=gray!10!white,
  coltitle=black,
  fonttitle=\bfseries\sffamily\small,
  title={Unit III $\bullet$ Card 2},
  boxrule=1.2pt,
  arc=2mm,
  outer arc=2mm,
  boxsep=1.5mm,
  top=1mm,
  bottom=1mm,
  left=1.5mm,
  right=1.5mm,
  before skip=2.5mm,
  after skip=2.5mm
]
\textbf{Q:} What are the key differences between Classic Bluetooth and Bluetooth Low Energy (BLE)?
\tcbline
\textbf{Ans:} \begin{itemize}
  \item \textbf{Classic Bluetooth:} High-throughput continuous data stream, high power consumption (~1W), latency up to 6s, max 7 slaves.
  \item \textbf{BLE:} Low duty cycle, ultra-low power consumption (~0.01W), latency $<$3ms, unlimited mesh topology nodes.
\end{itemize}
\end{tcolorbox}


\begin{tcolorbox}[
  colback=white,
  colframe=UnitColorThree,
  colbacktitle=gray!10!white,
  coltitle=black,
  fonttitle=\bfseries\sffamily\small,
  title={Unit III $\bullet$ Card 3},
  boxrule=1.2pt,
  arc=2mm,
  outer arc=2mm,
  boxsep=1.5mm,
  top=1mm,
  bottom=1mm,
  left=1.5mm,
  right=1.5mm,
  before skip=2.5mm,
  after skip=2.5mm
]
\textbf{Q:} Differentiate between Static and Dynamic IP addressing in IoT.
\tcbline
\textbf{Ans:} \textbf{Static IP} is manually configured and permanent; ideal for central network gateways or brokers. \textbf{Dynamic IP} is automatically assigned for a temporary lease duration by a DHCP server; ideal for edge sensor nodes.
\end{tcolorbox}


\begin{tcolorbox}[
  colback=white,
  colframe=UnitColorThree,
  colbacktitle=gray!10!white,
  coltitle=black,
  fonttitle=\bfseries\sffamily\small,
  title={Unit III $\bullet$ Card 4},
  boxrule=1.2pt,
  arc=2mm,
  outer arc=2mm,
  boxsep=1.5mm,
  top=1mm,
  bottom=1mm,
  left=1.5mm,
  right=1.5mm,
  before skip=2.5mm,
  after skip=2.5mm
]
\textbf{Q:} Explain the DORA handshake process in DHCP.
\tcbline
\textbf{Ans:} \textbf{DORA} stands for:
\begin{enumerate}
  \item \textbf{D}iscover: Client broadcasts to find a DHCP server.
  \item \textbf{O}ffer: Server proposes an IP address and network config.
  \item \textbf{R}equest: Client requests to lease the proposed IP.
  \item \textbf{A}cknowledge (ACK): Server commits the lease to the client.
\end{enumerate}
\end{tcolorbox}


\begin{tcolorbox}[
  colback=white,
  colframe=UnitColorThree,
  colbacktitle=gray!10!white,
  coltitle=black,
  fonttitle=\bfseries\sffamily\small,
  title={Unit III $\bullet$ Card 5},
  boxrule=1.2pt,
  arc=2mm,
  outer arc=2mm,
  boxsep=1.5mm,
  top=1mm,
  bottom=1mm,
  left=1.5mm,
  right=1.5mm,
  before skip=2.5mm,
  after skip=2.5mm
]
\textbf{Q:} Provide a brief example of DHCP dynamically configuring an IoT node on boot.
\tcbline
\textbf{Ans:} A smart thermostat boots up, broadcasts a `DHCPDISCOVER` packet over Wi-Fi. The home router replies with a `DHCPOFFER` proposing `192.168.1.15`. The thermostat accepts via `DHCPREQUEST`, and the router registers the thermostat's MAC address and confirms with a `DHCPACK`.
\end{tcolorbox}


\begin{tcolorbox}[
  colback=white,
  colframe=UnitColorThree,
  colbacktitle=gray!10!white,
  coltitle=black,
  fonttitle=\bfseries\sffamily\small,
  title={Unit III $\bullet$ Card 6},
  boxrule=1.2pt,
  arc=2mm,
  outer arc=2mm,
  boxsep=1.5mm,
  top=1mm,
  bottom=1mm,
  left=1.5mm,
  right=1.5mm,
  before skip=2.5mm,
  after skip=2.5mm
]
\textbf{Q:} Differentiate between TCP and UDP in terms of connection style, reliability, header overhead, and speed.
\tcbline
\textbf{Ans:} \begin{itemize}
  \item \textbf{TCP:} Connection-oriented (3-way handshake), reliable (ACKs/retransmissions), 20-byte min header, slower speed.
  \item \textbf{UDP:} Connectionless, best-effort (no delivery guarantee), 8-byte header, faster speed (low latency).
\end{itemize}
\end{tcolorbox}


\begin{tcolorbox}[
  colback=white,
  colframe=UnitColorThree,
  colbacktitle=gray!10!white,
  coltitle=black,
  fonttitle=\bfseries\sffamily\small,
  title={Unit III $\bullet$ Card 7},
  boxrule=1.2pt,
  arc=2mm,
  outer arc=2mm,
  boxsep=1.5mm,
  top=1mm,
  bottom=1mm,
  left=1.5mm,
  right=1.5mm,
  before skip=2.5mm,
  after skip=2.5mm
]
\textbf{Q:} Why is UDP preferred over TCP in resource-constrained IoT settings?
\tcbline
\textbf{Ans:} \begin{enumerate}
  \item \textbf{Power Conservation:} No connection handshakes/teardowns, saving radio uptime battery.
  \item \textbf{Low Overhead:} 8-byte header (vs TCP's 20-byte) saves precious network bandwidth.
  \item \textbf{Real-time Performance:} Dropped packets are ignored instead of causing lag via retransmission queues.
\end{enumerate}
\end{tcolorbox}


\begin{tcolorbox}[
  colback=white,
  colframe=UnitColorThree,
  colbacktitle=gray!10!white,
  coltitle=black,
  fonttitle=\bfseries\sffamily\small,
  title={Unit III $\bullet$ Card 8},
  boxrule=1.2pt,
  arc=2mm,
  outer arc=2mm,
  boxsep=1.5mm,
  top=1mm,
  bottom=1mm,
  left=1.5mm,
  right=1.5mm,
  before skip=2.5mm,
  after skip=2.5mm
]
\textbf{Q:} List five primary benefits of using IPv6 instead of IPv4 in IoT networks.
\tcbline
\textbf{Ans:} \begin{enumerate}
  \item Massive address space ($3.4 \times 10^{38}$ IPs).
  \item SLAAC (Stateless Address Autoconfiguration).
  \item Simplified, fixed-length 40-byte headers.
  \item Native IPsec security.
  \item Optimized multicasting (replaces wasteful broadcasts).
\end{enumerate}
\end{tcolorbox}


\begin{tcolorbox}[
  colback=white,
  colframe=UnitColorThree,
  colbacktitle=gray!10!white,
  coltitle=black,
  fonttitle=\bfseries\sffamily\small,
  title={Unit III $\bullet$ Card 9},
  boxrule=1.2pt,
  arc=2mm,
  outer arc=2mm,
  boxsep=1.5mm,
  top=1mm,
  bottom=1mm,
  left=1.5mm,
  right=1.5mm,
  before skip=2.5mm,
  after skip=2.5mm
]
\textbf{Q:} What is the length and structure of an IPv6 address?
\tcbline
\textbf{Ans:} \begin{itemize}
  \item \textbf{Length:} 128 bits (16 bytes).
  \item \textbf{Structure:} Written as 8 groups of 4 hexadecimal digits separated by colons (e.g., `2001:0db8:85a3:0000:0000:8a2e:0370:7334`). It contains a 64-bit Network Prefix and a 64-bit Interface Identifier.
\end{itemize}
\end{tcolorbox}


\begin{tcolorbox}[
  colback=white,
  colframe=UnitColorThree,
  colbacktitle=gray!10!white,
  coltitle=black,
  fonttitle=\bfseries\sffamily\small,
  title={Unit III $\bullet$ Card 10},
  boxrule=1.2pt,
  arc=2mm,
  outer arc=2mm,
  boxsep=1.5mm,
  top=1mm,
  bottom=1mm,
  left=1.5mm,
  right=1.5mm,
  before skip=2.5mm,
  after skip=2.5mm
]
\textbf{Q:} What is a MAC address, its standard length, and its role in a network?
\tcbline
\textbf{Ans:} A \textbf{MAC address} is a unique physical hardware identifier assigned to a network interface controller (NIC) by the manufacturer. It is \textbf{48 bits} (6 bytes) long, operates at the Data Link Layer (Layer 2), and uniquely identifies a physical node on a local network segment.
\end{tcolorbox}


\begin{tcolorbox}[
  colback=white,
  colframe=UnitColorThree,
  colbacktitle=gray!10!white,
  coltitle=black,
  fonttitle=\bfseries\sffamily\small,
  title={Unit III $\bullet$ Card 11},
  boxrule=1.2pt,
  arc=2mm,
  outer arc=2mm,
  boxsep=1.5mm,
  top=1mm,
  bottom=1mm,
  left=1.5mm,
  right=1.5mm,
  before skip=2.5mm,
  after skip=2.5mm
]
\textbf{Q:} Describe the internal bit structure of a 48-bit MAC address.
\tcbline
\textbf{Ans:} It is split into two halves:
\begin{enumerate}
  \item \textbf{OUI (Organizationally Unique Identifier):} First 24 bits, identifying the manufacturer (assigned by IEEE). Contains U/L (Universal/Local) and I/G (Individual/Group) bits.
  \item \textbf{NIC Specific ID:} Last 24 bits, uniquely assigned to the individual chip by the manufacturer.
\end{enumerate}
\end{tcolorbox}


\begin{tcolorbox}[
  colback=white,
  colframe=UnitColorThree,
  colbacktitle=gray!10!white,
  coltitle=black,
  fonttitle=\bfseries\sffamily\small,
  title={Unit III $\bullet$ Card 12},
  boxrule=1.2pt,
  arc=2mm,
  outer arc=2mm,
  boxsep=1.5mm,
  top=1mm,
  bottom=1mm,
  left=1.5mm,
  right=1.5mm,
  before skip=2.5mm,
  after skip=2.5mm
]
\textbf{Q:} What are network ports, and what are their standard classifications?
\tcbline
\textbf{Ans:} Ports are \textbf{16-bit logical endpoints} (values 0–65535) in the OS that route incoming network traffic to specific software processes.  
\textit{Classifications:} Well-Known Ports (0–1023), Registered Ports (1024–49151), and Dynamic/Private Ports (49152–65535).
\end{tcolorbox}


\begin{tcolorbox}[
  colback=white,
  colframe=UnitColorThree,
  colbacktitle=gray!10!white,
  coltitle=black,
  fonttitle=\bfseries\sffamily\small,
  title={Unit III $\bullet$ Card 13},
  boxrule=1.2pt,
  arc=2mm,
  outer arc=2mm,
  boxsep=1.5mm,
  top=1mm,
  bottom=1mm,
  left=1.5mm,
  right=1.5mm,
  before skip=2.5mm,
  after skip=2.5mm
]
\textbf{Q:} Explain the working principle, strengths, and limitations of HTTP in IoT.
\tcbline
\textbf{Ans:} \begin{itemize}
  \item \textbf{Principle:} Synchronous Request-Response protocol over TCP.
  \item \textbf{Strengths:} Ubiquitous, easy integration with web databases, secure (HTTPS/TLS).
  \item \textbf{Limitations:} High text-header overhead ($>$100s of bytes), synchronous pull-only model (forces wasteful client polling), heavy memory footprint.
\end{itemize}
\end{tcolorbox}


\begin{tcolorbox}[
  colback=white,
  colframe=UnitColorThree,
  colbacktitle=gray!10!white,
  coltitle=black,
  fonttitle=\bfseries\sffamily\small,
  title={Unit III $\bullet$ Card 14},
  boxrule=1.2pt,
  arc=2mm,
  outer arc=2mm,
  boxsep=1.5mm,
  top=1mm,
  bottom=1mm,
  left=1.5mm,
  right=1.5mm,
  before skip=2.5mm,
  after skip=2.5mm
]
\textbf{Q:} Explain the working principle, strengths, and limitations of CoAP in IoT.
\tcbline
\textbf{Ans:} \begin{itemize}
  \item \textbf{Principle:} RESTful Client-Server protocol running over UDP, using compact binary headers (4 bytes min) and supporting Observation mode.
  \item \textbf{Strengths:} Tiny packet overhead, asynchronous pushes, optimized for resource-constrained microcontrollers.
  \item \textbf{Limitations:} Transport-layer unreliability (requires custom application-layer error handling), NAT traversal issues with UDP.
\end{itemize}
\end{tcolorbox}


\begin{tcolorbox}[
  colback=white,
  colframe=UnitColorThree,
  colbacktitle=gray!10!white,
  coltitle=black,
  fonttitle=\bfseries\sffamily\small,
  title={Unit III $\bullet$ Card 15},
  boxrule=1.2pt,
  arc=2mm,
  outer arc=2mm,
  boxsep=1.5mm,
  top=1mm,
  bottom=1mm,
  left=1.5mm,
  right=1.5mm,
  before skip=2.5mm,
  after skip=2.5mm
]
\textbf{Q:} Explain the working principle, strengths, and limitations of MQTT in IoT.
\tcbline
\textbf{Ans:} \begin{itemize}
  \item \textbf{Principle:} Broker-based, event-driven Publish/Subscribe messaging protocol over TCP.
  \item \textbf{Strengths:} Tiny binary headers (2 bytes min), decoupled one-to-many communication, three QoS levels, Last Will and Testament (LWT) support.
  \item \textbf{Limitations:} Centralized broker acts as a single point of failure, TCP connection keep-alives drain battery, no native security (needs TLS).
\end{itemize}
\end{tcolorbox}


\begin{tcolorbox}[
  colback=white,
  colframe=UnitColorThree,
  colbacktitle=gray!10!white,
  coltitle=black,
  fonttitle=\bfseries\sffamily\small,
  title={Unit III $\bullet$ Card 16},
  boxrule=1.2pt,
  arc=2mm,
  outer arc=2mm,
  boxsep=1.5mm,
  top=1mm,
  bottom=1mm,
  left=1.5mm,
  right=1.5mm,
  before skip=2.5mm,
  after skip=2.5mm
]
\textbf{Q:} Summarize the transport protocol, architecture model, and minimum header size of HTTP, CoAP, and MQTT.
\tcbline
\textbf{Ans:} \begin{itemize}
  \item \textbf{HTTP:} TCP | Client-Server (Req-Resp) | 100+ bytes (ASCII text headers).
  \item \textbf{CoAP:} UDP | Client-Server (RESTful + Observe) | 4 bytes (Binary header).
  \item \textbf{MQTT:} TCP | Broker-based (Publish-Subscribe) | 2 bytes (Binary header).
\end{itemize}
\end{tcolorbox}

\section*{Unit IV: Chapter 04 Flashcards}

\begin{tcolorbox}[
  colback=white,
  colframe=UnitColorFour,
  colbacktitle=gray!10!white,
  coltitle=black,
  fonttitle=\bfseries\sffamily\small,
  title={Unit IV $\bullet$ Card 1},
  boxrule=1.2pt,
  arc=2mm,
  outer arc=2mm,
  boxsep=1.5mm,
  top=1mm,
  bottom=1mm,
  left=1.5mm,
  right=1.5mm,
  before skip=2.5mm,
  after skip=2.5mm
]
\textbf{Q:} What is the primary goal of prototyping in IoT?
\tcbline
\textbf{Ans:} To create a tangible, low-risk representation of an IoT concept to validate design assumptions, test hardware/software integration, check user interactions, and resolve bugs before mass manufacturing.
\end{tcolorbox}


\begin{tcolorbox}[
  colback=white,
  colframe=UnitColorFour,
  colbacktitle=gray!10!white,
  coltitle=black,
  fonttitle=\bfseries\sffamily\small,
  title={Unit IV $\bullet$ Card 2},
  boxrule=1.2pt,
  arc=2mm,
  outer arc=2mm,
  boxsep=1.5mm,
  top=1mm,
  bottom=1mm,
  left=1.5mm,
  right=1.5mm,
  before skip=2.5mm,
  after skip=2.5mm
]
\textbf{Q:} What is the core trade-off represented by the "Cost vs. Ease of Prototyping" curve?
\tcbline
\textbf{Ans:} As development progresses from initial concept to mass manufacture, the cost of implementing a design change increases exponentially, while the ease of making changes and design flexibility drops to zero.
\end{tcolorbox}


\begin{tcolorbox}[
  colback=white,
  colframe=UnitColorFour,
  colbacktitle=gray!10!white,
  coltitle=black,
  fonttitle=\bfseries\sffamily\small,
  title={Unit IV $\bullet$ Card 3},
  boxrule=1.2pt,
  arc=2mm,
  outer arc=2mm,
  boxsep=1.5mm,
  top=1mm,
  bottom=1mm,
  left=1.5mm,
  right=1.5mm,
  before skip=2.5mm,
  after skip=2.5mm
]
\textbf{Q:} List the five phases of the prototyping curve in order.
\tcbline
\textbf{Ans:} \begin{enumerate}
  \item \textbf{Sketching \& Paper Mockups} (near-zero cost, infinite ease).
  \item \textbf{Breadboarding / Sandbox} (low cost, high ease).
  \item \textbf{Engineering Prototype} (medium cost, moderate ease).
  \item \textbf{Pilot Run} (high cost, low ease).
  \item \textbf{Mass Production} (extreme cost, zero ease).
\end{enumerate}
\end{tcolorbox}


\begin{tcolorbox}[
  colback=white,
  colframe=UnitColorFour,
  colbacktitle=gray!10!white,
  coltitle=black,
  fonttitle=\bfseries\sffamily\small,
  title={Unit IV $\bullet$ Card 4},
  boxrule=1.2pt,
  arc=2mm,
  outer arc=2mm,
  boxsep=1.5mm,
  top=1mm,
  bottom=1mm,
  left=1.5mm,
  right=1.5mm,
  before skip=2.5mm,
  after skip=2.5mm
]
\textbf{Q:} In what four ways does prototyping bridge the gap to final production?
\tcbline
\textbf{Ans:} \begin{enumerate}
  \item Identifies electrical noise or battery issues.
  \item Tests firmware connection and logical loop stability.
  \item Verifies mechanical fit (PCB inside casing).
  \item Gathers early user experience (UX) feedback.
\end{enumerate}
\end{tcolorbox}


\begin{tcolorbox}[
  colback=white,
  colframe=UnitColorFour,
  colbacktitle=gray!10!white,
  coltitle=black,
  fonttitle=\bfseries\sffamily\small,
  title={Unit IV $\bullet$ Card 5},
  boxrule=1.2pt,
  arc=2mm,
  outer arc=2mm,
  boxsep=1.5mm,
  top=1mm,
  bottom=1mm,
  left=1.5mm,
  right=1.5mm,
  before skip=2.5mm,
  after skip=2.5mm
]
\textbf{Q:} Compare Open-Source and Closed-Source prototyping paradigms in IoT.
\tcbline
\textbf{Ans:} \begin{itemize}
  \item \textbf{Open Source:} Low upfront cost, community-driven support, fully custom toolchains, massive shared libraries (e.g., Arduino, ESP32).
  \item \textbf{Closed Source:} High license costs, direct vendor SLA support, locked proprietary interfaces, pre-certified hardware modules (e.g., Electric Imp).
\end{itemize}
\end{tcolorbox}


\begin{tcolorbox}[
  colback=white,
  colframe=UnitColorFour,
  colbacktitle=gray!10!white,
  coltitle=black,
  fonttitle=\bfseries\sffamily\small,
  title={Unit IV $\bullet$ Card 6},
  boxrule=1.2pt,
  arc=2mm,
  outer arc=2mm,
  boxsep=1.5mm,
  top=1mm,
  bottom=1mm,
  left=1.5mm,
  right=1.5mm,
  before skip=2.5mm,
  after skip=2.5mm
]
\textbf{Q:} What are the primary advantages and disadvantages of using Open Source in IoT prototyping?
\tcbline
\textbf{Ans:} \begin{itemize}
  \item \textit{Advantages:} Zero vendor lock-in, free pre-existing driver libraries, and low entry cost.
  \item \textit{Disadvantages:} No official SLA/technical support guarantees, risk of copyleft license compliance, and inconsistent documentation.
\end{itemize}
\end{tcolorbox}


\begin{tcolorbox}[
  colback=white,
  colframe=UnitColorFour,
  colbacktitle=gray!10!white,
  coltitle=black,
  fonttitle=\bfseries\sffamily\small,
  title={Unit IV $\bullet$ Card 7},
  boxrule=1.2pt,
  arc=2mm,
  outer arc=2mm,
  boxsep=1.5mm,
  top=1mm,
  bottom=1mm,
  left=1.5mm,
  right=1.5mm,
  before skip=2.5mm,
  after skip=2.5mm
]
\textbf{Q:} What is the role of "Sketching" in the early phases of IoT development?
\tcbline
\textbf{Ans:} Sketching refers to creating low-fidelity, non-functional models (pencil layouts, cardboard/foam mockups, Figma wireframes) to quickly test device scale, ergonomics, and button/LED positioning without writing any code.
\end{tcolorbox}


\begin{tcolorbox}[
  colback=white,
  colframe=UnitColorFour,
  colbacktitle=gray!10!white,
  coltitle=black,
  fonttitle=\bfseries\sffamily\small,
  title={Unit IV $\bullet$ Card 8},
  boxrule=1.2pt,
  arc=2mm,
  outer arc=2mm,
  boxsep=1.5mm,
  top=1mm,
  bottom=1mm,
  left=1.5mm,
  right=1.5mm,
  before skip=2.5mm,
  after skip=2.5mm
]
\textbf{Q:} How do "Familiarity" and "Community" accelerate the prototyping process?
\tcbline
\textbf{Ans:} \begin{itemize}
  \item \textbf{Familiarity:} Using known platforms (e.g., Python, Arduino IDE) bypasses the learning curve, shifting focus to solving the core product problems.
  \item \textbf{Community:} Tapping into open-source forums and code repositories allows developers to find immediate fixes for bugs, saving months of manual driver development.
\end{itemize}
\end{tcolorbox}

\section*{Unit V: Chapter 05 Flashcards}

\begin{tcolorbox}[
  colback=white,
  colframe=UnitColorFive,
  colbacktitle=gray!10!white,
  coltitle=black,
  fonttitle=\bfseries\sffamily\small,
  title={Unit V $\bullet$ Card 1},
  boxrule=1.2pt,
  arc=2mm,
  outer arc=2mm,
  boxsep=1.5mm,
  top=1mm,
  bottom=1mm,
  left=1.5mm,
  right=1.5mm,
  before skip=2.5mm,
  after skip=2.5mm
]
\textbf{Q:} Explain the working principle of 3D Printing (FDM) in IoT prototyping.
\tcbline
\textbf{Ans:} Fused Deposition Modeling (FDM) melts thermoplastic filament and extrudes it through a nozzle, depositing it layer-by-layer on a build plate based on a sliced 3D CAD model to form a solid enclosure.
\end{tcolorbox}


\begin{tcolorbox}[
  colback=white,
  colframe=UnitColorFive,
  colbacktitle=gray!10!white,
  coltitle=black,
  fonttitle=\bfseries\sffamily\small,
  title={Unit V $\bullet$ Card 2},
  boxrule=1.2pt,
  arc=2mm,
  outer arc=2mm,
  boxsep=1.5mm,
  top=1mm,
  bottom=1mm,
  left=1.5mm,
  right=1.5mm,
  before skip=2.5mm,
  after skip=2.5mm
]
\textbf{Q:} What are the main limitations of 3D Printing for physical prototyping?
\tcbline
\textbf{Ans:} Slow build speed (takes hours), anisotropic structural weakness (parts easily snap along layer lines), and rough surface finishes showing layer lines.
\end{tcolorbox}


\begin{tcolorbox}[
  colback=white,
  colframe=UnitColorFive,
  colbacktitle=gray!10!white,
  coltitle=black,
  fonttitle=\bfseries\sffamily\small,
  title={Unit V $\bullet$ Card 3},
  boxrule=1.2pt,
  arc=2mm,
  outer arc=2mm,
  boxsep=1.5mm,
  top=1mm,
  bottom=1mm,
  left=1.5mm,
  right=1.5mm,
  before skip=2.5mm,
  after skip=2.5mm
]
\textbf{Q:} What is the primary dimensional constraint of Laser Cutting and how do designers bypass it?
\tcbline
\textbf{Ans:} It is restricted to 2D flat profiles. Designers bypass this by cutting flat sheets with integrated tabs or finger joints, which are then assembled/slotted together to form 3D boxes.
\end{tcolorbox}


\begin{tcolorbox}[
  colback=white,
  colframe=UnitColorFive,
  colbacktitle=gray!10!white,
  coltitle=black,
  fonttitle=\bfseries\sffamily\small,
  title={Unit V $\bullet$ Card 4},
  boxrule=1.2pt,
  arc=2mm,
  outer arc=2mm,
  boxsep=1.5mm,
  top=1mm,
  bottom=1mm,
  left=1.5mm,
  right=1.5mm,
  before skip=2.5mm,
  after skip=2.5mm
]
\textbf{Q:} Why must PVC (Polyvinyl Chloride) never be laser cut?
\tcbline
\textbf{Ans:} Laser cutting PVC releases toxic and highly corrosive chlorine gas, which presents extreme safety hazards to operators and corrodes the metal components of the laser machine.
\end{tcolorbox}


\begin{tcolorbox}[
  colback=white,
  colframe=UnitColorFive,
  colbacktitle=gray!10!white,
  coltitle=black,
  fonttitle=\bfseries\sffamily\small,
  title={Unit V $\bullet$ Card 5},
  boxrule=1.2pt,
  arc=2mm,
  outer arc=2mm,
  boxsep=1.5mm,
  top=1mm,
  bottom=1mm,
  left=1.5mm,
  right=1.5mm,
  before skip=2.5mm,
  after skip=2.5mm
]
\textbf{Q:} Detail the working principle, main advantage, and main disadvantage of CNC Milling.
\tcbline
\textbf{Ans:} CNC milling is a subtractive process that carves shapes out of a solid block using high-speed rotating cutting tools.  
\begin{itemize}
  \item \textit{Advantage:} High mechanical strength and precise tolerances.
  \item \textit{Disadvantage:} High material waste and complex setup (generating G-code).
\end{itemize}
\end{tcolorbox}


\begin{tcolorbox}[
  colback=white,
  colframe=UnitColorFive,
  colbacktitle=gray!10!white,
  coltitle=black,
  fonttitle=\bfseries\sffamily\small,
  title={Unit V $\bullet$ Card 6},
  boxrule=1.2pt,
  arc=2mm,
  outer arc=2mm,
  boxsep=1.5mm,
  top=1mm,
  bottom=1mm,
  left=1.5mm,
  right=1.5mm,
  before skip=2.5mm,
  after skip=2.5mm
]
\textbf{Q:} What are the advantages and limitations of repurposing off-the-shelf enclosures?
\tcbline
\textbf{Ans:} \begin{itemize}
  \item \textit{Advantages:} Pre-fabricated water resistance (IP-rated project boxes), low cost, and instant availability.
  \item \textit{Limitations:} Restricted to pre-existing dimensions and hard to scale or replicate identically.
\end{itemize}
\end{tcolorbox}


\begin{tcolorbox}[
  colback=white,
  colframe=UnitColorFive,
  colbacktitle=gray!10!white,
  coltitle=black,
  fonttitle=\bfseries\sffamily\small,
  title={Unit V $\bullet$ Card 7},
  boxrule=1.2pt,
  arc=2mm,
  outer arc=2mm,
  boxsep=1.5mm,
  top=1mm,
  bottom=1mm,
  left=1.5mm,
  right=1.5mm,
  before skip=2.5mm,
  after skip=2.5mm
]
\textbf{Q:} What is the function of an antenna in an IoT device?
\tcbline
\textbf{Ans:} It acts as a transducer, converting electrical signals from the RF transceiver chip into electromagnetic radio waves for transmission, and capturing incoming waves to convert them back to electrical signals.
\end{tcolorbox}


\begin{tcolorbox}[
  colback=white,
  colframe=UnitColorFive,
  colbacktitle=gray!10!white,
  coltitle=black,
  fonttitle=\bfseries\sffamily\small,
  title={Unit V $\bullet$ Card 8},
  boxrule=1.2pt,
  arc=2mm,
  outer arc=2mm,
  boxsep=1.5mm,
  top=1mm,
  bottom=1mm,
  left=1.5mm,
  right=1.5mm,
  before skip=2.5mm,
  after skip=2.5mm
]
\textbf{Q:} Compare PCB Trace and External Whip antennas in terms of unit cost and range.
\tcbline
\textbf{Ans:} \begin{itemize}
  \item \textbf{PCB Trace:} Zero unit cost (fabricated as copper traces on the board), moderate range.
  \item \textbf{External Whip:} High unit cost (requires coax cables and SMA connectors), maximum range and high gain.
\end{itemize}
\end{tcolorbox}


\begin{tcolorbox}[
  colback=white,
  colframe=UnitColorFive,
  colbacktitle=gray!10!white,
  coltitle=black,
  fonttitle=\bfseries\sffamily\small,
  title={Unit V $\bullet$ Card 9},
  boxrule=1.2pt,
  arc=2mm,
  outer arc=2mm,
  boxsep=1.5mm,
  top=1mm,
  bottom=1mm,
  left=1.5mm,
  right=1.5mm,
  before skip=2.5mm,
  after skip=2.5mm
]
\textbf{Q:} How does a metallic enclosure affect an internal antenna, and how is it resolved?
\tcbline
\textbf{Ans:} A metal enclosure acts as a Faraday cage, blocking electromagnetic RF signals. This is resolved by mounting an external antenna outside the enclosure connected to the PCB via an SMA feedthrough.
\end{tcolorbox}


\begin{tcolorbox}[
  colback=white,
  colframe=UnitColorFive,
  colbacktitle=gray!10!white,
  coltitle=black,
  fonttitle=\bfseries\sffamily\small,
  title={Unit V $\bullet$ Card 10},
  boxrule=1.2pt,
  arc=2mm,
  outer arc=2mm,
  boxsep=1.5mm,
  top=1mm,
  bottom=1mm,
  left=1.5mm,
  right=1.5mm,
  before skip=2.5mm,
  after skip=2.5mm
]
\textbf{Q:} What is Ground Plane Clearance and why is it necessary for chip antennas?
\tcbline
\textbf{Ans:} It is a design requirement stating that no copper ground planes, traces, or vias should run on any layer of the PCB directly underneath or adjacent to the antenna. This is necessary to prevent detuning the antenna and blocking the RF signal.
\end{tcolorbox}

\section*{Unit VI: Chapter 06 Flashcards}

\begin{tcolorbox}[
  colback=white,
  colframe=UnitColorSix,
  colbacktitle=gray!10!white,
  coltitle=black,
  fonttitle=\bfseries\sffamily\small,
  title={Unit VI $\bullet$ Card 1},
  boxrule=1.2pt,
  arc=2mm,
  outer arc=2mm,
  boxsep=1.5mm,
  top=1mm,
  bottom=1mm,
  left=1.5mm,
  right=1.5mm,
  before skip=2.5mm,
  after skip=2.5mm
]
\textbf{Q:} What is the fundamental architectural difference between Arduino and Raspberry Pi?
\tcbline
\textbf{Ans:} Arduino is a \textbf{microcontroller} platform (8-bit ATmega, bare-metal firmware, no OS, instant boot). Raspberry Pi is a \textbf{microprocessor} platform (64-bit ARM Cortex-A72 SoC, runs a full Linux OS, 15–30 second boot).
\end{tcolorbox}


\begin{tcolorbox}[
  colback=white,
  colframe=UnitColorSix,
  colbacktitle=gray!10!white,
  coltitle=black,
  fonttitle=\bfseries\sffamily\small,
  title={Unit VI $\bullet$ Card 2},
  boxrule=1.2pt,
  arc=2mm,
  outer arc=2mm,
  boxsep=1.5mm,
  top=1mm,
  bottom=1mm,
  left=1.5mm,
  right=1.5mm,
  before skip=2.5mm,
  after skip=2.5mm
]
\textbf{Q:} Compare Arduino and Raspberry Pi across RAM, connectivity, and power consumption.
\tcbline
\textbf{Ans:} \begin{itemize}
  \item \textbf{Arduino:} 2 KB SRAM, no built-in Wi-Fi/BLE (needs shields), ultra-low power (~50 mA).
  \item \textbf{Raspberry Pi:} 1–8 GB LPDDR4, built-in Wi-Fi/BLE/Ethernet, higher power (~600 mA–1.2 A).
\end{itemize}
\end{tcolorbox}


\begin{tcolorbox}[
  colback=white,
  colframe=UnitColorSix,
  colbacktitle=gray!10!white,
  coltitle=black,
  fonttitle=\bfseries\sffamily\small,
  title={Unit VI $\bullet$ Card 3},
  boxrule=1.2pt,
  arc=2mm,
  outer arc=2mm,
  boxsep=1.5mm,
  top=1mm,
  bottom=1mm,
  left=1.5mm,
  right=1.5mm,
  before skip=2.5mm,
  after skip=2.5mm
]
\textbf{Q:} When should you use Arduino over Raspberry Pi for an IoT project?
\tcbline
\textbf{Ans:} Use Arduino for single-task, deterministic, real-time control loops requiring microsecond-level timing (e.g., reading a sensor every 100 ms and toggling a relay), where low power consumption and instant boot are critical.
\end{tcolorbox}


\begin{tcolorbox}[
  colback=white,
  colframe=UnitColorSix,
  colbacktitle=gray!10!white,
  coltitle=black,
  fonttitle=\bfseries\sffamily\small,
  title={Unit VI $\bullet$ Card 4},
  boxrule=1.2pt,
  arc=2mm,
  outer arc=2mm,
  boxsep=1.5mm,
  top=1mm,
  bottom=1mm,
  left=1.5mm,
  right=1.5mm,
  before skip=2.5mm,
  after skip=2.5mm
]
\textbf{Q:} List five strengths of the Raspberry Pi for IoT prototyping.
\tcbline
\textbf{Ans:} \begin{enumerate}
  \item Full Linux OS support (thousands of packages).
  \item High processing power (quad-core ARM, 1.5 GHz).
  \item Built-in Wi-Fi, BLE, and Gigabit Ethernet.
  \item Massive community and documentation.
  \item Versatile 40-pin GPIO header (I2C, SPI, UART).
\end{enumerate}
\end{tcolorbox}


\begin{tcolorbox}[
  colback=white,
  colframe=UnitColorSix,
  colbacktitle=gray!10!white,
  coltitle=black,
  fonttitle=\bfseries\sffamily\small,
  title={Unit VI $\bullet$ Card 5},
  boxrule=1.2pt,
  arc=2mm,
  outer arc=2mm,
  boxsep=1.5mm,
  top=1mm,
  bottom=1mm,
  left=1.5mm,
  right=1.5mm,
  before skip=2.5mm,
  after skip=2.5mm
]
\textbf{Q:} List five limitations of the Raspberry Pi for IoT.
\tcbline
\textbf{Ans:} \begin{enumerate}
  \item No built-in ADC (needs external MCP3008).
  \item High power consumption (~1.2 A under load).
  \item Non-real-time OS (cannot guarantee microsecond timing).
  \item SD card filesystem corruption risk on power loss.
  \item Thermal throttling under sustained CPU load.
\end{enumerate}
\end{tcolorbox}


\begin{tcolorbox}[
  colback=white,
  colframe=UnitColorSix,
  colbacktitle=gray!10!white,
  coltitle=black,
  fonttitle=\bfseries\sffamily\small,
  title={Unit VI $\bullet$ Card 6},
  boxrule=1.2pt,
  arc=2mm,
  outer arc=2mm,
  boxsep=1.5mm,
  top=1mm,
  bottom=1mm,
  left=1.5mm,
  right=1.5mm,
  before skip=2.5mm,
  after skip=2.5mm
]
\textbf{Q:} Name four operating systems supported by the Raspberry Pi.
\tcbline
\textbf{Ans:} Raspberry Pi OS (Debian-based, official), Ubuntu Server/Desktop (ARM64), Windows 10 IoT Core (headless UWP apps), and RISC OS (lightweight, fast-boot niche OS).
\end{tcolorbox}


\begin{tcolorbox}[
  colback=white,
  colframe=UnitColorSix,
  colbacktitle=gray!10!white,
  coltitle=black,
  fonttitle=\bfseries\sffamily\small,
  title={Unit VI $\bullet$ Card 7},
  boxrule=1.2pt,
  arc=2mm,
  outer arc=2mm,
  boxsep=1.5mm,
  top=1mm,
  bottom=1mm,
  left=1.5mm,
  right=1.5mm,
  before skip=2.5mm,
  after skip=2.5mm
]
\textbf{Q:} Define GPIO pins and name the three communication buses available on the Raspberry Pi GPIO header.
\tcbline
\textbf{Ans:} \textbf{GPIO} pins are programmable digital pins configurable as Input or Output. The three communication buses are:  
\begin{enumerate}
  \item \textbf{I2C} (SDA + SCL, shared multi-device bus).
  \item \textbf{SPI} (MOSI, MISO, SCLK, CS, high-speed full-duplex).
  \item \textbf{UART} (TX + RX, point-to-point serial).
\end{enumerate}
\end{tcolorbox}


\begin{tcolorbox}[
  colback=white,
  colframe=UnitColorSix,
  colbacktitle=gray!10!white,
  coltitle=black,
  fonttitle=\bfseries\sffamily\small,
  title={Unit VI $\bullet$ Card 8},
  boxrule=1.2pt,
  arc=2mm,
  outer arc=2mm,
  boxsep=1.5mm,
  top=1mm,
  bottom=1mm,
  left=1.5mm,
  right=1.5mm,
  before skip=2.5mm,
  after skip=2.5mm
]
\textbf{Q:} State five critical safety guidelines for Raspberry Pi GPIO pins.
\tcbline
\textbf{Ans:} \begin{enumerate}
  \item Never exceed 3.3V on any GPIO pin.
  \item Never draw more than 16 mA per pin.
  \item Never short 5V to 3.3V or GND.
  \item Always use current-limiting resistors for LEDs.
  \item Never hot-plug wires while the Pi is powered on.
\end{enumerate}
\end{tcolorbox}


\begin{tcolorbox}[
  colback=white,
  colframe=UnitColorSix,
  colbacktitle=gray!10!white,
  coltitle=black,
  fonttitle=\bfseries\sffamily\small,
  title={Unit VI $\bullet$ Card 9},
  boxrule=1.2pt,
  arc=2mm,
  outer arc=2mm,
  boxsep=1.5mm,
  top=1mm,
  bottom=1mm,
  left=1.5mm,
  right=1.5mm,
  before skip=2.5mm,
  after skip=2.5mm
]
\textbf{Q:} What are the two unique hardware advantages of the BeagleBone Black over the Raspberry Pi?
\tcbline
\textbf{Ans:} \begin{enumerate}
  \item \textbf{Built-in 12-bit ADC} (7 analog input channels at 1.8V max), unlike the Pi which has no native ADC.
  \item \textbf{PRU (Programmable Real-Time Units):} Two independent 200 MHz co-processors for deterministic, microsecond-level I/O alongside Linux.
\end{enumerate}
\end{tcolorbox}


\begin{tcolorbox}[
  colback=white,
  colframe=UnitColorSix,
  colbacktitle=gray!10!white,
  coltitle=black,
  fonttitle=\bfseries\sffamily\small,
  title={Unit VI $\bullet$ Card 10},
  boxrule=1.2pt,
  arc=2mm,
  outer arc=2mm,
  boxsep=1.5mm,
  top=1mm,
  bottom=1mm,
  left=1.5mm,
  right=1.5mm,
  before skip=2.5mm,
  after skip=2.5mm
]
\textbf{Q:} Explain the architecture and primary limitation of the Electric Imp platform.
\tcbline
\textbf{Ans:} Electric Imp splits device logic into Device Code (runs on the physical module) and Cloud Code (runs on Electric Imp's managed servers), using the proprietary Squirrel language. \textbf{Primary limitation:} Complete vendor lock-in — if the vendor shuts down its cloud, all devices become non-functional.
\end{tcolorbox}


\begin{tcolorbox}[
  colback=white,
  colframe=UnitColorSix,
  colbacktitle=gray!10!white,
  coltitle=black,
  fonttitle=\bfseries\sffamily\small,
  title={Unit VI $\bullet$ Card 11},
  boxrule=1.2pt,
  arc=2mm,
  outer arc=2mm,
  boxsep=1.5mm,
  top=1mm,
  bottom=1mm,
  left=1.5mm,
  right=1.5mm,
  before skip=2.5mm,
  after skip=2.5mm
]
\textbf{Q:} Why are smartphones and tablets useful as IoT prototyping platforms, and what is their biggest limitation?
\tcbline
\textbf{Ans:} They have rich built-in sensor arrays (accelerometer, GPS, camera, microphone) and connectivity (Wi-Fi, BLE, NFC, Cellular), eliminating the need for separate breakout boards. \textbf{Biggest limitation:} Not designed for always-on, unattended operation (battery drain, OS updates, thermal shutdowns).
\end{tcolorbox}


\begin{tcolorbox}[
  colback=white,
  colframe=UnitColorSix,
  colbacktitle=gray!10!white,
  coltitle=black,
  fonttitle=\bfseries\sffamily\small,
  title={Unit VI $\bullet$ Card 12},
  boxrule=1.2pt,
  arc=2mm,
  outer arc=2mm,
  boxsep=1.5mm,
  top=1mm,
  bottom=1mm,
  left=1.5mm,
  right=1.5mm,
  before skip=2.5mm,
  after skip=2.5mm
]
\textbf{Q:} What is a Plug Computer and what is its primary IoT advantage?
\tcbline
\textbf{Ans:} A Plug Computer is a small, wall-socket-mounted computing device (e.g., SheevaPlug) that runs Linux and draws power directly from mains electricity. \textbf{Primary advantage:} Always-on by design, making it ideal for 24/7 home automation hubs or local MQTT brokers without battery concerns.
\end{tcolorbox}

\section*{Unit VII: Chapter 07 Flashcards}

\begin{tcolorbox}[
  colback=white,
  colframe=UnitColorSeven,
  colbacktitle=gray!10!white,
  coltitle=black,
  fonttitle=\bfseries\sffamily\small,
  title={Unit VII $\bullet$ Card 1},
  boxrule=1.2pt,
  arc=2mm,
  outer arc=2mm,
  boxsep=1.5mm,
  top=1mm,
  bottom=1mm,
  left=1.5mm,
  right=1.5mm,
  before skip=2.5mm,
  after skip=2.5mm
]
\textbf{Q:} What is an API in the context of IoT?
\tcbline
\textbf{Ans:} An \textbf{API (Application Programming Interface)} is a set of defined rules and protocols allowing different software applications to communicate. In IoT, APIs bridge physical devices (sensors/actuators) and online services (cloud platforms, databases, third-party apps).
\end{tcolorbox}


\begin{tcolorbox}[
  colback=white,
  colframe=UnitColorSeven,
  colbacktitle=gray!10!white,
  coltitle=black,
  fonttitle=\bfseries\sffamily\small,
  title={Unit VII $\bullet$ Card 2},
  boxrule=1.2pt,
  arc=2mm,
  outer arc=2mm,
  boxsep=1.5mm,
  top=1mm,
  bottom=1mm,
  left=1.5mm,
  right=1.5mm,
  before skip=2.5mm,
  after skip=2.5mm
]
\textbf{Q:} Distinguish between consuming an existing API and writing a new API for IoT.
\tcbline
\textbf{Ans:} \begin{itemize}
  \item \textbf{Consuming:} The device acts as a client, sending HTTP requests to a third-party cloud API to fetch/push data (e.g., querying OpenWeatherMap).
  \item \textbf{Writing:} The device acts as a server, exposing its own RESTful endpoints so external apps can read sensors or control actuators (e.g., `GET /api/sensors/soil-moisture`).
\end{itemize}
\end{tcolorbox}


\begin{tcolorbox}[
  colback=white,
  colframe=UnitColorSeven,
  colbacktitle=gray!10!white,
  coltitle=black,
  fonttitle=\bfseries\sffamily\small,
  title={Unit VII $\bullet$ Card 3},
  boxrule=1.2pt,
  arc=2mm,
  outer arc=2mm,
  boxsep=1.5mm,
  top=1mm,
  bottom=1mm,
  left=1.5mm,
  right=1.5mm,
  before skip=2.5mm,
  after skip=2.5mm
]
\textbf{Q:} Define the role of Cloud Computing in IoT and list its four primary functions.
\tcbline
\textbf{Ans:} Cloud computing provides remote, scalable infrastructure for IoT. Four functions:  
\begin{enumerate}
  \item \textbf{Data Storage} (unlimited scalable archives).
  \item \textbf{Data Processing \& Analytics} (ML training, stream processing).
  \item \textbf{Device Management} (registration, OTA updates, device twins).
  \item \textbf{Application Hosting} (dashboards, mobile backends, notification services).
\end{enumerate}
\end{tcolorbox}


\begin{tcolorbox}[
  colback=white,
  colframe=UnitColorSeven,
  colbacktitle=gray!10!white,
  coltitle=black,
  fonttitle=\bfseries\sffamily\small,
  title={Unit VII $\bullet$ Card 4},
  boxrule=1.2pt,
  arc=2mm,
  outer arc=2mm,
  boxsep=1.5mm,
  top=1mm,
  bottom=1mm,
  left=1.5mm,
  right=1.5mm,
  before skip=2.5mm,
  after skip=2.5mm
]
\textbf{Q:} Explain the role of Cloud Computing in a Smart Grid architecture.
\tcbline
\textbf{Ans:} Smart Meters transmit real-time consumption data to the cloud. The cloud performs demand forecasting, load balancing, and peak shaving. It sends control commands to grid substations and provides consumer-facing apps for real-time billing and usage alerts.
\end{tcolorbox}


\begin{tcolorbox}[
  colback=white,
  colframe=UnitColorSeven,
  colbacktitle=gray!10!white,
  coltitle=black,
  fonttitle=\bfseries\sffamily\small,
  title={Unit VII $\bullet$ Card 5},
  boxrule=1.2pt,
  arc=2mm,
  outer arc=2mm,
  boxsep=1.5mm,
  top=1mm,
  bottom=1mm,
  left=1.5mm,
  right=1.5mm,
  before skip=2.5mm,
  after skip=2.5mm
]
\textbf{Q:} Name three advantages and three disadvantages of cloud computing in Smart Grids.
\tcbline
\textbf{Ans:} \begin{itemize}
  \item \textit{Advantages:} Centralized processing of millions of meters, elastic scaling during peak demand, and long-term historical trend analysis.
  \item \textit{Disadvantages:} Latency too slow for protection relay decisions, granular data reveals private household behaviors, and WAN outages disconnect all meters.
\end{itemize}
\end{tcolorbox}


\begin{tcolorbox}[
  colback=white,
  colframe=UnitColorSeven,
  colbacktitle=gray!10!white,
  coltitle=black,
  fonttitle=\bfseries\sffamily\small,
  title={Unit VII $\bullet$ Card 6},
  boxrule=1.2pt,
  arc=2mm,
  outer arc=2mm,
  boxsep=1.5mm,
  top=1mm,
  bottom=1mm,
  left=1.5mm,
  right=1.5mm,
  before skip=2.5mm,
  after skip=2.5mm
]
\textbf{Q:} Describe the four levels of Big Data Analytics in IoT.
\tcbline
\textbf{Ans:} \begin{enumerate}
  \item \textbf{Descriptive} ("What happened?") — Dashboards and summaries.
  \item \textbf{Diagnostic} ("Why did it happen?") — Root cause analysis.
  \item \textbf{Predictive} ("What will happen?") — ML forecasting.
  \item \textbf{Prescriptive} ("What should we do?") — Automated optimal actions.
\end{enumerate}
\end{tcolorbox}


\begin{tcolorbox}[
  colback=white,
  colframe=UnitColorSeven,
  colbacktitle=gray!10!white,
  coltitle=black,
  fonttitle=\bfseries\sffamily\small,
  title={Unit VII $\bullet$ Card 7},
  boxrule=1.2pt,
  arc=2mm,
  outer arc=2mm,
  boxsep=1.5mm,
  top=1mm,
  bottom=1mm,
  left=1.5mm,
  right=1.5mm,
  before skip=2.5mm,
  after skip=2.5mm
]
\textbf{Q:} Give a concrete IoT example for Predictive and Prescriptive analytics.
\tcbline
\textbf{Ans:} \begin{itemize}
  \item \textbf{Predictive:} A motor vibration model forecasts ball bearing failure in 14 days based on 6-month frequency patterns.
  \item \textbf{Prescriptive:} The system auto-schedules a maintenance work order, reserves spare parts, and assigns a technician without human intervention.
\end{itemize}
\end{tcolorbox}


\begin{tcolorbox}[
  colback=white,
  colframe=UnitColorSeven,
  colbacktitle=gray!10!white,
  coltitle=black,
  fonttitle=\bfseries\sffamily\small,
  title={Unit VII $\bullet$ Card 8},
  boxrule=1.2pt,
  arc=2mm,
  outer arc=2mm,
  boxsep=1.5mm,
  top=1mm,
  bottom=1mm,
  left=1.5mm,
  right=1.5mm,
  before skip=2.5mm,
  after skip=2.5mm
]
\textbf{Q:} What are Real-Time Reactions in IoT? Give an example.
\tcbline
\textbf{Ans:} The ability of an IoT system to process an incoming event and trigger an immediate, automated response. \textbf{Example:} A motion sensor detects an intruder; within 200 ms the system triggers a siren, locks smart doors, turns on floodlights, and pushes a phone notification.
\end{tcolorbox}


\begin{tcolorbox}[
  colback=white,
  colframe=UnitColorSeven,
  colbacktitle=gray!10!white,
  coltitle=black,
  fonttitle=\bfseries\sffamily\small,
  title={Unit VII $\bullet$ Card 9},
  boxrule=1.2pt,
  arc=2mm,
  outer arc=2mm,
  boxsep=1.5mm,
  top=1mm,
  bottom=1mm,
  left=1.5mm,
  right=1.5mm,
  before skip=2.5mm,
  after skip=2.5mm
]
\textbf{Q:} How do WebSockets differ from standard HTTP for IoT communication?
\tcbline
\textbf{Ans:} HTTP is request-response only (client must poll). WebSockets upgrade a single TCP connection into a \textbf{persistent, full-duplex, bidirectional} channel where both client and server can push messages at any time without re-establishing a connection.
\end{tcolorbox}


\begin{tcolorbox}[
  colback=white,
  colframe=UnitColorSeven,
  colbacktitle=gray!10!white,
  coltitle=black,
  fonttitle=\bfseries\sffamily\small,
  title={Unit VII $\bullet$ Card 10},
  boxrule=1.2pt,
  arc=2mm,
  outer arc=2mm,
  boxsep=1.5mm,
  top=1mm,
  bottom=1mm,
  left=1.5mm,
  right=1.5mm,
  before skip=2.5mm,
  after skip=2.5mm
]
\textbf{Q:} Compare HTTP Polling, WebSockets, and MQTT for real-time IoT.
\tcbline
\textbf{Ans:} \begin{itemize}
  \item \textbf{HTTP Polling:} High latency, client-initiated only, high overhead per request.
  \item \textbf{WebSockets:} Very low latency, full-duplex bidirectional, ideal for browser dashboards.
  \item \textbf{MQTT:} Low latency, broker-mediated pub/sub push, minimal 2-byte header, ideal for massive device networks.
\end{itemize}
\end{tcolorbox}

\section*{Unit VIII: Chapter 08 Flashcards}

\begin{tcolorbox}[
  colback=white,
  colframe=UnitColorEight,
  colbacktitle=gray!10!white,
  coltitle=black,
  fonttitle=\bfseries\sffamily\small,
  title={Unit VIII $\bullet$ Card 1},
  boxrule=1.2pt,
  arc=2mm,
  outer arc=2mm,
  boxsep=1.5mm,
  top=1mm,
  bottom=1mm,
  left=1.5mm,
  right=1.5mm,
  before skip=2.5mm,
  after skip=2.5mm
]
\textbf{Q:} Why should dynamic memory allocation be avoided on microcontrollers?
\tcbline
\textbf{Ans:} Calling `malloc()`/`free()` at runtime causes \textbf{heap fragmentation} — free memory splits into non-contiguous blocks too small for new allocations. Microcontrollers lack virtual memory, so fragmentation leads to crashes. Use fixed-size, statically allocated buffers instead.
\end{tcolorbox}


\begin{tcolorbox}[
  colback=white,
  colframe=UnitColorEight,
  colbacktitle=gray!10!white,
  coltitle=black,
  fonttitle=\bfseries\sffamily\small,
  title={Unit VIII $\bullet$ Card 2},
  boxrule=1.2pt,
  arc=2mm,
  outer arc=2mm,
  boxsep=1.5mm,
  top=1mm,
  bottom=1mm,
  left=1.5mm,
  right=1.5mm,
  before skip=2.5mm,
  after skip=2.5mm
]
\textbf{Q:} What is the difference between blocking delays and non-blocking timer-based state machines?
\tcbline
\textbf{Ans:} `delay()` halts the entire CPU, wasting cycles and blocking other tasks. Non-blocking approaches use `millis()` or hardware timer interrupts to schedule work, allowing the CPU to handle other events while waiting.
\end{tcolorbox}


\begin{tcolorbox}[
  colback=white,
  colframe=UnitColorEight,
  colbacktitle=gray!10!white,
  coltitle=black,
  fonttitle=\bfseries\sffamily\small,
  title={Unit VIII $\bullet$ Card 3},
  boxrule=1.2pt,
  arc=2mm,
  outer arc=2mm,
  boxsep=1.5mm,
  top=1mm,
  bottom=1mm,
  left=1.5mm,
  right=1.5mm,
  before skip=2.5mm,
  after skip=2.5mm
]
\textbf{Q:} Name three techniques for optimizing battery life in embedded IoT firmware.
\tcbline
\textbf{Ans:} \begin{enumerate}
  \item Leverage deep sleep modes (power down CPU/radio, wake on timer/GPIO).
  \item Reduce radio duty cycle (short, infrequent burst transmissions).
  \item Lower clock speed when idle via Dynamic Voltage and Frequency Scaling (DVFS).
\end{enumerate}
\end{tcolorbox}


\begin{tcolorbox}[
  colback=white,
  colframe=UnitColorEight,
  colbacktitle=gray!10!white,
  coltitle=black,
  fonttitle=\bfseries\sffamily\small,
  title={Unit VIII $\bullet$ Card 4},
  boxrule=1.2pt,
  arc=2mm,
  outer arc=2mm,
  boxsep=1.5mm,
  top=1mm,
  bottom=1mm,
  left=1.5mm,
  right=1.5mm,
  before skip=2.5mm,
  after skip=2.5mm
]
\textbf{Q:} What should you audit before adopting a third-party library for an embedded project?
\tcbline
\textbf{Ans:} Compile the library and check its \textbf{Flash and RAM consumption} using the compiler's memory usage report. Avoid bloated frameworks designed for desktops; prefer lightweight, single-purpose embedded libraries.
\end{tcolorbox}


\begin{tcolorbox}[
  colback=white,
  colframe=UnitColorEight,
  colbacktitle=gray!10!white,
  coltitle=black,
  fonttitle=\bfseries\sffamily\small,
  title={Unit VIII $\bullet$ Card 5},
  boxrule=1.2pt,
  arc=2mm,
  outer arc=2mm,
  boxsep=1.5mm,
  top=1mm,
  bottom=1mm,
  left=1.5mm,
  right=1.5mm,
  before skip=2.5mm,
  after skip=2.5mm
]
\textbf{Q:} Compare Serial Print debugging and JTAG/SWD hardware debugging.
\tcbline
\textbf{Ans:} \begin{itemize}
  \item \textbf{Serial Print:} Simple, no special hardware. But it modifies code timing (masks race conditions) and consumes UART and RAM.
  \item \textbf{JTAG/SWD:} Non-intrusive, allows breakpoints and real-time register inspection. But requires a physical debug probe and complex setup.
\end{itemize}
\end{tcolorbox}


\begin{tcolorbox}[
  colback=white,
  colframe=UnitColorEight,
  colbacktitle=gray!10!white,
  coltitle=black,
  fonttitle=\bfseries\sffamily\small,
  title={Unit VIII $\bullet$ Card 6},
  boxrule=1.2pt,
  arc=2mm,
  outer arc=2mm,
  boxsep=1.5mm,
  top=1mm,
  bottom=1mm,
  left=1.5mm,
  right=1.5mm,
  before skip=2.5mm,
  after skip=2.5mm
]
\textbf{Q:} Name five common funding channels for IoT startups.
\tcbline
\textbf{Ans:} \begin{enumerate}
  \item Bootstrapping (self-funding).
  \item Crowdfunding (Kickstarter/Indiegogo).
  \item Angel investors (seed equity).
  \item Venture capital (Series A/B/C).
  \item Government grants and hardware accelerators (SBIR, HAX, Techstars).
\end{enumerate}
\end{tcolorbox}


\begin{tcolorbox}[
  colback=white,
  colframe=UnitColorEight,
  colbacktitle=gray!10!white,
  coltitle=black,
  fonttitle=\bfseries\sffamily\small,
  title={Unit VIII $\bullet$ Card 7},
  boxrule=1.2pt,
  arc=2mm,
  outer arc=2mm,
  boxsep=1.5mm,
  top=1mm,
  bottom=1mm,
  left=1.5mm,
  right=1.5mm,
  before skip=2.5mm,
  after skip=2.5mm
]
\textbf{Q:} What is the Business Model Canvas and how many building blocks does it contain?
\tcbline
\textbf{Ans:} The BMC (Alexander Osterwalder) is a single-page strategic tool mapping \textbf{9 building blocks}: Customer Segments, Value Proposition, Channels, Customer Relationships, Revenue Streams, Key Resources, Key Activities, Key Partnerships, and Cost Structure.
\end{tcolorbox}


\begin{tcolorbox}[
  colback=white,
  colframe=UnitColorEight,
  colbacktitle=gray!10!white,
  coltitle=black,
  fonttitle=\bfseries\sffamily\small,
  title={Unit VIII $\bullet$ Card 8},
  boxrule=1.2pt,
  arc=2mm,
  outer arc=2mm,
  boxsep=1.5mm,
  top=1mm,
  bottom=1mm,
  left=1.5mm,
  right=1.5mm,
  before skip=2.5mm,
  after skip=2.5mm
]
\textbf{Q:} Describe the Build-Measure-Learn loop in Lean Startup methodology applied to IoT.
\tcbline
\textbf{Ans:} \begin{enumerate}
  \item \textbf{Build:} Create a Minimum Viable Product (MVP) — simplest functional prototype.
  \item \textbf{Measure:} Deploy to early adopters; collect usage data and error rates.
  \item \textbf{Learn:} Analyze data to validate or invalidate the core hypothesis. Then \textbf{Persevere} or \textbf{Pivot}.
\end{enumerate}
\end{tcolorbox}


\begin{tcolorbox}[
  colback=white,
  colframe=UnitColorEight,
  colbacktitle=gray!10!white,
  coltitle=black,
  fonttitle=\bfseries\sffamily\small,
  title={Unit VIII $\bullet$ Card 9},
  boxrule=1.2pt,
  arc=2mm,
  outer arc=2mm,
  boxsep=1.5mm,
  top=1mm,
  bottom=1mm,
  left=1.5mm,
  right=1.5mm,
  before skip=2.5mm,
  after skip=2.5mm
]
\textbf{Q:} What does "Pivot or Persevere" mean in the context of an IoT startup?
\tcbline
\textbf{Ans:} After each Build-Measure-Learn cycle, the startup decides: \textbf{Persevere} (data validates the hypothesis, continue the current direction) or \textbf{Pivot} (data invalidates the hypothesis, change the customer segment, value proposition, or technology platform).
\end{tcolorbox}


\begin{tcolorbox}[
  colback=white,
  colframe=UnitColorEight,
  colbacktitle=gray!10!white,
  coltitle=black,
  fonttitle=\bfseries\sffamily\small,
  title={Unit VIII $\bullet$ Card 10},
  boxrule=1.2pt,
  arc=2mm,
  outer arc=2mm,
  boxsep=1.5mm,
  top=1mm,
  bottom=1mm,
  left=1.5mm,
  right=1.5mm,
  before skip=2.5mm,
  after skip=2.5mm
]
\textbf{Q:} State two advantages and two disadvantages of applying Lean Startup to IoT hardware ventures.
\tcbline
\textbf{Ans:} \begin{itemize}
  \item \textit{Advantages:} Reduces financial risk before expensive mass production; accelerates time-to-market by focusing on validated features.
  \item \textit{Disadvantages:} Hardware MVPs are slower and costlier to iterate than software MVPs; early adopters may not represent the broader market.
\end{itemize}
\end{tcolorbox}

\section*{Unit IX: Chapter 09 Flashcards}

\begin{tcolorbox}[
  colback=white,
  colframe=UnitColorNine,
  colbacktitle=gray!10!white,
  coltitle=black,
  fonttitle=\bfseries\sffamily\small,
  title={Unit IX $\bullet$ Card 1},
  boxrule=1.2pt,
  arc=2mm,
  outer arc=2mm,
  boxsep=1.5mm,
  top=1mm,
  bottom=1mm,
  left=1.5mm,
  right=1.5mm,
  before skip=2.5mm,
  after skip=2.5mm
]
\textbf{Q:} List the seven steps in the PCB design-to-manufacturing pipeline in order.
\tcbline
\textbf{Ans:} \begin{enumerate}
  \item Schematic Capture.
  \item PCB Layout (Board Design).
  \item Design Rule Check (DRC).
  \item Gerber File Generation.
  \item PCB Fabrication.
  \item Component Assembly (SMT/THT).
  \item Testing and Quality Assurance.
\end{enumerate}
\end{tcolorbox}


\begin{tcolorbox}[
  colback=white,
  colframe=UnitColorNine,
  colbacktitle=gray!10!white,
  coltitle=black,
  fonttitle=\bfseries\sffamily\small,
  title={Unit IX $\bullet$ Card 2},
  boxrule=1.2pt,
  arc=2mm,
  outer arc=2mm,
  boxsep=1.5mm,
  top=1mm,
  bottom=1mm,
  left=1.5mm,
  right=1.5mm,
  before skip=2.5mm,
  after skip=2.5mm
]
\textbf{Q:} What is a Gerber file and why is it important?
\tcbline
\textbf{Ans:} A Gerber file is the universal, industry-standard format for PCB manufacturing. A complete set includes a separate file for each copper layer, solder mask layer, silkscreen layer, and drill file. It is the deliverable sent to the fabrication house.
\end{tcolorbox}


\begin{tcolorbox}[
  colback=white,
  colframe=UnitColorNine,
  colbacktitle=gray!10!white,
  coltitle=black,
  fonttitle=\bfseries\sffamily\small,
  title={Unit IX $\bullet$ Card 3},
  boxrule=1.2pt,
  arc=2mm,
  outer arc=2mm,
  boxsep=1.5mm,
  top=1mm,
  bottom=1mm,
  left=1.5mm,
  right=1.5mm,
  before skip=2.5mm,
  after skip=2.5mm
]
\textbf{Q:} Differentiate between SMT and Through-Hole assembly techniques.
\tcbline
\textbf{Ans:} \begin{itemize}
  \item \textbf{SMT (Surface Mount):} Pick-and-place machine positions tiny components on solder paste pads; reflow oven melts and bonds them. Fast and automated.
  \item \textbf{Through-Hole (THT):} Larger components are inserted through drilled holes and soldered via wave soldering or manual soldering. Stronger mechanical bonds.
\end{itemize}
\end{tcolorbox}


\begin{tcolorbox}[
  colback=white,
  colframe=UnitColorNine,
  colbacktitle=gray!10!white,
  coltitle=black,
  fonttitle=\bfseries\sffamily\small,
  title={Unit IX $\bullet$ Card 4},
  boxrule=1.2pt,
  arc=2mm,
  outer arc=2mm,
  boxsep=1.5mm,
  top=1mm,
  bottom=1mm,
  left=1.5mm,
  right=1.5mm,
  before skip=2.5mm,
  after skip=2.5mm
]
\textbf{Q:} Name the four testing methods used after PCB assembly.
\tcbline
\textbf{Ans:} \begin{enumerate}
  \item \textbf{AOI} (Automated Optical Inspection) — camera scan for defects.
  \item \textbf{ICT} (In-Circuit Testing) — electrical probe for continuity.
  \item \textbf{Functional Test} — power on, flash firmware, verify all systems.
  \item \textbf{Burn-In Test} — sustained stress operation for 24–72 hours.
\end{enumerate}
\end{tcolorbox}


\begin{tcolorbox}[
  colback=white,
  colframe=UnitColorNine,
  colbacktitle=gray!10!white,
  coltitle=black,
  fonttitle=\bfseries\sffamily\small,
  title={Unit IX $\bullet$ Card 5},
  boxrule=1.2pt,
  arc=2mm,
  outer arc=2mm,
  boxsep=1.5mm,
  top=1mm,
  bottom=1mm,
  left=1.5mm,
  right=1.5mm,
  before skip=2.5mm,
  after skip=2.5mm
]
\textbf{Q:} What are the five categories of testing an IoT system must undergo before mass manufacture?
\tcbline
\textbf{Ans:} \begin{enumerate}
  \item Hardware Testing (ESD, thermal, drop/vibration, EMI).
  \item Software/Firmware Testing (unit, integration, regression, OTA).
  \item Connectivity Testing (Wi-Fi range, BLE pairing, MQTT QoS).
  \item Security Testing (pen testing, firmware analysis, TLS validation).
  \item User Acceptance Testing (onboarding, mobile app UX, field trials).
\end{enumerate}
\end{tcolorbox}


\begin{tcolorbox}[
  colback=white,
  colframe=UnitColorNine,
  colbacktitle=gray!10!white,
  coltitle=black,
  fonttitle=\bfseries\sffamily\small,
  title={Unit IX $\bullet$ Card 6},
  boxrule=1.2pt,
  arc=2mm,
  outer arc=2mm,
  boxsep=1.5mm,
  top=1mm,
  bottom=1mm,
  left=1.5mm,
  right=1.5mm,
  before skip=2.5mm,
  after skip=2.5mm
]
\textbf{Q:} Explain ESD testing and thermal cycling in the context of IoT hardware validation.
\tcbline
\textbf{Ans:} \begin{itemize}
  \item \textbf{ESD Testing:} Subjects the device to high-voltage static discharge pulses (IEC 61000-4-2) to ensure survival during real-world handling.
  \item \textbf{Thermal Cycling:} Repeatedly cycles between extreme temperatures (e.g., -20°C to +70°C) to test solder joint reliability and component tolerance.
\end{itemize}
\end{tcolorbox}


\begin{tcolorbox}[
  colback=white,
  colframe=UnitColorNine,
  colbacktitle=gray!10!white,
  coltitle=black,
  fonttitle=\bfseries\sffamily\small,
  title={Unit IX $\bullet$ Card 7},
  boxrule=1.2pt,
  arc=2mm,
  outer arc=2mm,
  boxsep=1.5mm,
  top=1mm,
  bottom=1mm,
  left=1.5mm,
  right=1.5mm,
  before skip=2.5mm,
  after skip=2.5mm
]
\textbf{Q:} What is the purpose of firmware binary analysis in IoT security testing?
\tcbline
\textbf{Ans:} The firmware binary is extracted and decompiled to check for hardcoded credentials, API keys, unencrypted sensitive data, and insecure coding patterns that an attacker could exploit.
\end{tcolorbox}


\begin{tcolorbox}[
  colback=white,
  colframe=UnitColorNine,
  colbacktitle=gray!10!white,
  coltitle=black,
  fonttitle=\bfseries\sffamily\small,
  title={Unit IX $\bullet$ Card 8},
  boxrule=1.2pt,
  arc=2mm,
  outer arc=2mm,
  boxsep=1.5mm,
  top=1mm,
  bottom=1mm,
  left=1.5mm,
  right=1.5mm,
  before skip=2.5mm,
  after skip=2.5mm
]
\textbf{Q:} What is injection molding and what are its key advantages and disadvantages for IoT enclosures?
\tcbline
\textbf{Ans:} Molten plastic is injected into a precision steel mold at high pressure; after cooling, the part is ejected.  
\begin{itemize}
  \item \textit{Advantage:} Produces thousands of identical, smooth parts per day at very low per-unit cost.
  \item \textit{Disadvantage:} Extremely high upfront mold tooling cost (\$10K–\$100K+); any design change requires a new mold.
\end{itemize}
\end{tcolorbox}


\begin{tcolorbox}[
  colback=white,
  colframe=UnitColorNine,
  colbacktitle=gray!10!white,
  coltitle=black,
  fonttitle=\bfseries\sffamily\small,
  title={Unit IX $\bullet$ Card 9},
  boxrule=1.2pt,
  arc=2mm,
  outer arc=2mm,
  boxsep=1.5mm,
  top=1mm,
  bottom=1mm,
  left=1.5mm,
  right=1.5mm,
  before skip=2.5mm,
  after skip=2.5mm
]
\textbf{Q:} Name four regulatory certifications required before selling an IoT product commercially.
\tcbline
\textbf{Ans:} \begin{enumerate}
  \item \textbf{FCC Part 15} (US — RF emissions limits).
  \item \textbf{CE Marking} (EU — health, safety, environment conformity).
  \item \textbf{RoHS} (EU/Global — restriction of hazardous substances).
  \item \textbf{UL / IEC 62368-1} (Global — electrical safety).
\end{enumerate}
\end{tcolorbox}


\begin{tcolorbox}[
  colback=white,
  colframe=UnitColorNine,
  colbacktitle=gray!10!white,
  coltitle=black,
  fonttitle=\bfseries\sffamily\small,
  title={Unit IX $\bullet$ Card 10},
  boxrule=1.2pt,
  arc=2mm,
  outer arc=2mm,
  boxsep=1.5mm,
  top=1mm,
  bottom=1mm,
  left=1.5mm,
  right=1.5mm,
  before skip=2.5mm,
  after skip=2.5mm
]
\textbf{Q:} What are the four key aspects of scaling up software when moving from prototype to mass production?
\tcbline
\textbf{Ans:} \begin{enumerate}
  \item Auto-scaling cloud infrastructure (AWS ECS, Kubernetes).
  \item Centralized device fleet management and bulk OTA updates.
  \item Distributed NoSQL/time-series databases for high-velocity data.
  \item CI/CD pipelines for automated firmware and cloud releases.
\end{enumerate}
\end{tcolorbox}

\section*{Unit X: Chapter 10 Flashcards}

\begin{tcolorbox}[
  colback=white,
  colframe=UnitColorTen,
  colbacktitle=gray!10!white,
  coltitle=black,
  fonttitle=\bfseries\sffamily\small,
  title={Unit X $\bullet$ Card 1},
  boxrule=1.2pt,
  arc=2mm,
  outer arc=2mm,
  boxsep=1.5mm,
  top=1mm,
  bottom=1mm,
  left=1.5mm,
  right=1.5mm,
  before skip=2.5mm,
  after skip=2.5mm
]
\textbf{Q:} Define a Smart City and explain the role IoT plays in its implementation.
\tcbline
\textbf{Ans:} A Smart City is an urban area that uses ICT and IoT sensors to gather data, manage resources, and run municipal utilities efficiently. IoT nodes act as the sensory network (detecting traffic, light, and trash fill-levels) to optimize power, logistics, and service delivery.
\end{tcolorbox}


\begin{tcolorbox}[
  colback=white,
  colframe=UnitColorTen,
  colbacktitle=gray!10!white,
  coltitle=black,
  fonttitle=\bfseries\sffamily\small,
  title={Unit X $\bullet$ Card 2},
  boxrule=1.2pt,
  arc=2mm,
  outer arc=2mm,
  boxsep=1.5mm,
  top=1mm,
  bottom=1mm,
  left=1.5mm,
  right=1.5mm,
  before skip=2.5mm,
  after skip=2.5mm
]
\textbf{Q:} What is the working principle and benefits of Smart Lighting in a Smart City?
\tcbline
\textbf{Ans:} \begin{itemize}
  \item \textbf{Working Principle:} Streetlights use Light Dependent Resistors (LDRs) and motion sensors (PIR) to adjust light intensity based on time, ambient light, and activity.
  \item \textbf{Benefits:} Saves up to 80\% energy during low-traffic periods, and automatically reports bulb failures to a central maintenance panel.
\end{itemize}
\end{tcolorbox}


\begin{tcolorbox}[
  colback=white,
  colframe=UnitColorTen,
  colbacktitle=gray!10!white,
  coltitle=black,
  fonttitle=\bfseries\sffamily\small,
  title={Unit X $\bullet$ Card 3},
  boxrule=1.2pt,
  arc=2mm,
  outer arc=2mm,
  boxsep=1.5mm,
  top=1mm,
  bottom=1mm,
  left=1.5mm,
  right=1.5mm,
  before skip=2.5mm,
  after skip=2.5mm
]
\textbf{Q:} Explain how dynamic signal control and emergency preemption function in smart traffic management.
\tcbline
\textbf{Ans:} \begin{itemize}
  \item \textbf{Dynamic Signal Control:} Cameras/sensors monitor traffic volume at intersections and adjust signal phases in real-time to clear actual queues rather than using rigid timers.
  \item \textbf{Emergency Preemption:} Traffic signals track approaching emergency vehicles (ambulances/fire engines) via GPS and automatically switch to green to provide priority routing.
\end{itemize}
\end{tcolorbox}


\begin{tcolorbox}[
  colback=white,
  colframe=UnitColorTen,
  colbacktitle=gray!10!white,
  coltitle=black,
  fonttitle=\bfseries\sffamily\small,
  title={Unit X $\bullet$ Card 4},
  boxrule=1.2pt,
  arc=2mm,
  outer arc=2mm,
  boxsep=1.5mm,
  top=1mm,
  bottom=1mm,
  left=1.5mm,
  right=1.5mm,
  before skip=2.5mm,
  after skip=2.5mm
]
\textbf{Q:} Describe the working principle of ultrasonic waste tracking in smart cities.
\tcbline
\textbf{Ans:} Bins are fitted with ultrasonic sensors on their inner lids. The sensor measures the distance to the trash surface to track the fill level and sends this data via cellular modules (NB-IoT) to central routing software. This allows municipal trucks to skip empty bins and optimize collection routes.
\end{tcolorbox}


\begin{tcolorbox}[
  colback=white,
  colframe=UnitColorTen,
  colbacktitle=gray!10!white,
  coltitle=black,
  fonttitle=\bfseries\sffamily\small,
  title={Unit X $\bullet$ Card 5},
  boxrule=1.2pt,
  arc=2mm,
  outer arc=2mm,
  boxsep=1.5mm,
  top=1mm,
  bottom=1mm,
  left=1.5mm,
  right=1.5mm,
  before skip=2.5mm,
  after skip=2.5mm
]
\textbf{Q:} What is Vehicle-to-Everything (V2X) communication in Intelligent Transportation Systems (ITS)?
\tcbline
\textbf{Ans:} V2X communication is the transfer of information between a vehicle and any entity that may affect the vehicle (including other vehicles V2V, and roadside infrastructure V2I). It enables real-time public transit tracking, smart parking spot reservation, and vehicle collision warnings.
\end{tcolorbox}


\begin{tcolorbox}[
  colback=white,
  colframe=UnitColorTen,
  colbacktitle=gray!10!white,
  coltitle=black,
  fonttitle=\bfseries\sffamily\small,
  title={Unit X $\bullet$ Card 6},
  boxrule=1.2pt,
  arc=2mm,
  outer arc=2mm,
  boxsep=1.5mm,
  top=1mm,
  bottom=1mm,
  left=1.5mm,
  right=1.5mm,
  before skip=2.5mm,
  after skip=2.5mm
]
\textbf{Q:} List the four main pillars of ethical challenges in IoT.
\tcbline
\textbf{Ans:} \begin{enumerate}
  \item \textbf{Privacy} (surveillance creep, data profiling, consent ambiguity).
  \item \textbf{Control \& Ownership} (vendor lock-in, remote bricking, data custody).
  \item \textbf{Security Vulnerabilities} (botnets, physical safety risk, unpatched firmware).
  \item \textbf{Environmental Impact} (e-waste, planned obsolescence, high scale energy use).
\end{enumerate}
\end{tcolorbox}


\begin{tcolorbox}[
  colback=white,
  colframe=UnitColorTen,
  colbacktitle=gray!10!white,
  coltitle=black,
  fonttitle=\bfseries\sffamily\small,
  title={Unit X $\bullet$ Card 7},
  boxrule=1.2pt,
  arc=2mm,
  outer arc=2mm,
  boxsep=1.5mm,
  top=1mm,
  bottom=1mm,
  left=1.5mm,
  right=1.5mm,
  before skip=2.5mm,
  after skip=2.5mm
]
\textbf{Q:} Differentiate between "surveillance creep" and "data profiling" in the context of IoT privacy.
\tcbline
\textbf{Ans:} \begin{itemize}
  \item \textbf{Surveillance Creep:} The passive, constant monitoring of individuals in private/public spaces by ambient sensors without active consent.
  \item \textbf{Data Profiling:} Aggregating and analyzing multiple benign sensor data points (like appliance usage and TV logs) using machine learning to deduce highly personal schedules, habits, and details.
\end{itemize}
\end{tcolorbox}


\begin{tcolorbox}[
  colback=white,
  colframe=UnitColorTen,
  colbacktitle=gray!10!white,
  coltitle=black,
  fonttitle=\bfseries\sffamily\small,
  title={Unit X $\bullet$ Card 8},
  boxrule=1.2pt,
  arc=2mm,
  outer arc=2mm,
  boxsep=1.5mm,
  top=1mm,
  bottom=1mm,
  left=1.5mm,
  right=1.5mm,
  before skip=2.5mm,
  after skip=2.5mm
]
\textbf{Q:} What is remote bricking and how does it relate to vendor lock-in?
\tcbline
\textbf{Ans:} Remote bricking occurs when a vendor disables a device remotely or shuts down the cloud servers required for its operation, rendering functional hardware useless. This is a severe risk of vendor lock-in, where proprietary walled gardens prevent users from running their devices on open or alternative local platforms.
\end{tcolorbox}


\begin{tcolorbox}[
  colback=white,
  colframe=UnitColorTen,
  colbacktitle=gray!10!white,
  coltitle=black,
  fonttitle=\bfseries\sffamily\small,
  title={Unit X $\bullet$ Card 9},
  boxrule=1.2pt,
  arc=2mm,
  outer arc=2mm,
  boxsep=1.5mm,
  top=1mm,
  bottom=1mm,
  left=1.5mm,
  right=1.5mm,
  before skip=2.5mm,
  after skip=2.5mm
]
\textbf{Q:} How do insecure consumer IoT devices lead to Mirai-class botnets?
\tcbline
\textbf{Ans:} Millions of cheap IoT devices (IP cameras, baby monitors) are shipped with hardcoded default passwords and lack secure boot. Attackers scan the internet, compromise these devices using automated scripts, and recruit them into massive botnets to run Distributed Denial of Service (DDoS) attacks.
\end{tcolorbox}


\begin{tcolorbox}[
  colback=white,
  colframe=UnitColorTen,
  colbacktitle=gray!10!white,
  coltitle=black,
  fonttitle=\bfseries\sffamily\small,
  title={Unit X $\bullet$ Card 10},
  boxrule=1.2pt,
  arc=2mm,
  outer arc=2mm,
  boxsep=1.5mm,
  top=1mm,
  bottom=1mm,
  left=1.5mm,
  right=1.5mm,
  before skip=2.5mm,
  after skip=2.5mm
]
\textbf{Q:} Name four mitigation strategies that address the ethical challenges of IoT.
\tcbline
\textbf{Ans:} \begin{enumerate}
  \item \textbf{Privacy by Design:} Local edge processing and data minimization.
  \item \textbf{Open Standards:} Matter/Zigbee open APIs to prevent vendor lock-in.
  \item \textbf{Security Lifecycle:} Unique default passwords and mandatory security updates.
  \item \textbf{Circular Design:} Modular builds with replaceable batteries and recycling programs.
\end{enumerate}
\end{tcolorbox}

\end{document}